% concatenated from paper-lpprop.tex

\documentclass[a4paper, reqno, 11pt]{amsart}
\usepackage[utf8]{inputenc}
\usepackage{amsmath}
\usepackage{amssymb}
\usepackage{amsthm}
\usepackage[dvipsnames]{xcolor}
\usepackage{comment}


%\usepackage[backend=bibtex,style=alphabetic,sorting=ynt]{biblatex}
%\addbibresource{bibLPWE.bib}

\newtheorem{theorem}{Theorem}[section]
\newtheorem{corollary}[theorem]{Corollary}
\newtheorem{lemma}[theorem]{Lemma}
\newtheorem{claim}[theorem]{Claim}
\newtheorem{example}[theorem]{Example}
\newtheorem{proposition}[theorem]{Proposition}
\newtheorem{remark}[theorem]{Remark}
\theoremstyle{definition}
\newtheorem{definition}{Definition}[section]
%\theoremstyle{remark}
%\newtheorem*{remark}{Remark}

\DeclareMathOperator\supp{supp}

\newcommand{\philogm}{\varphi(\log\langle mx \rangle)}

\usepackage[luatex,colorlinks=true, linkcolor=WildStrawberry, citecolor=PineGreen]{hyperref}

\subjclass[MSC 2020]{35L05, 35R60}

\keywords{Wave equation, $L^p$ estimates, random data}

\title[A pathological set for a.s. properties of Gaussian measures]{A pathological set regarding the propagation of almost sure properties of Gaussian measures}
\author{Pablo Merino}
\address[P. Merino]{Basque Center for Applied Mathematics, 48009 Bilbao, Basque Country}
\date{}


\begin{document}
	\maketitle
	
	\begin{abstract}
	We provide a complementary result to the quasi-invariance of Gaussian measures supported on Sobolev spaces with high regularity under the dynamics of the three-dimensional periodic defocusing nonlinear wave equation from Gunaratnam-Oh-Tzvetkov-Weber (2022). Namely, given $p$ and $\sigma$ large enough, we prove the existence of dense sets of Sobolev spaces $W^{\sigma,p}(\mathbb{T}^3)$ which do not preserve the regularity $\sigma$ throughout the aforementioned dynamics. This is in sharp contrast with the propagation under the flow of almost sure properties.
	\end{abstract}
	
	%\tableofcontents
	
	%From having a dense set which does not preserve regularity, what kind of 'fine property'-consequences are there for the general flow?
	
	%What can we properly say, instead, w.r.t. the standard gaussian measure?
	

	\section{Introduction}\label{sec:intro}
	Let $\sigma,p \geq 1$. We consider the Cauchy problem associated to the defocusing cubic nonlinear wave equation on $\mathbb{T}^3 = (\mathbb{R}/\mathbb{Z})^3$:
		\begin{equation}\label{eq:cauchywave}
			\left\{
			\begin{array}{l}
				\partial_{tt} u - \Delta u + u^3 = 0,
				\\
				u(x,0) = 0,  \ \ \partial_tu(x,0) = g(x),
			\end{array}
			\right.
		\end{equation}
		where $u:\mathbb{T}^3 \times \mathbb{R} \rightarrow \mathbb{R}$ is the unknown function. The initial datum $g$ in \eqref{eq:cauchywave} is an element of the set
	\begin{align}\label{eq:family-S}
	S_{\lambda} = C^{\infty}(\mathbb{T}^3) + \{\delta f_{\lambda} : \delta > 0\},
	\end{align}
	where $f_{\lambda}$, for suitable $\lambda$ and $p$ (see Lemma \ref{lemma:flambda-wsigmap-per} below), are elements of $W^{\sigma-1,p}(\mathbb{T}^3)$ given by the Fourier coefficients (see \eqref{eq:flambdadef-1})
\begin{align}\label{eq:flambdadef}
\widehat{f_{\lambda}}(k) =  h(k) \langle k \rangle^{-\sigma+1} |k|^{-\lambda}, \quad k \in \mathbb{Z}^3,
\end{align}	
for $h$ a suitable function in $C^{\infty}(\mathbb{R}^3)$ which is zero near the origin (see \eqref{eq:def-h}). Given $d \in \mathbb{N}$, $\alpha \geq 0$ and $q \in [1,\infty]$, we denote by $W^{\alpha,q}(\mathbb{T}^d)$ (resp. $W^{\alpha,q}(\mathbb{R}^d)$) the set of distributions on $\mathbb{T}^d$ which weak derivatives of order up to $\alpha$ belong to $L^q(\mathbb{T}^d)$ (resp. $L^q(\mathbb{R}^d)$), endowed with the norm
        \begin{align*}
            \|f\|_{W^{\alpha,q}(\mathbb{T}^d)} := \|(1-\Delta)^{\frac{\alpha}{2}}f\|_{L^q(\mathbb{T}^d)} \quad (\text{resp. } \|f\|_{W^{\alpha,q}(\mathbb{R}^d)} := \|(1-\Delta)^{\frac{\alpha}{2}}f\|_{L^q(\mathbb{R}^d)} ).
        \end{align*} 
        For the same values of $\lambda$ and $p$, from the density of $C^{\infty}(\mathbb{T}^3)$ in $W^{\sigma-1,p}(\mathbb{T}^3)$ it follows that $S_{\lambda}$ is a dense subset of $W^{\sigma-1,p}(\mathbb{T}^3)$.  The homogeneous part of the solution of \eqref{eq:cauchywave} is given by
\begin{align}\label{eq:linearflow-flambda}
S(t)(0,g) = \frac{\sin(t|D|)}{|D|}g = \frac{e^{it|D|} - e^{-it|D|}}{2i|D|}g,
\end{align} 
where we denote $D = \frac{1}{i} \nabla$ (for $\nabla$ the gradient in $x$), $|D| = \sqrt{-\Delta}$ and $\langle D \rangle = (1-\Delta)^{\frac{1}{2}}$.

\hfill \break

In works \cite{Peral1980} and \cite{Miyachi1980}, Peral and Miyachi independently proved that, for fixed nonzero times, $L^p$ estimates on $\mathbb{R}^d$ exhibit a loss of derivatives under the linear wave evolution. Namely, they showed that, given $t \neq 0$ fixed,
\begin{align*}
\langle D \rangle^{-\sigma}e^{it|D|} \text{ is bounded on } L^p(\mathbb{R}^d) \Longleftrightarrow \sigma \geq (d-1) \Big| \frac{1}{p} - \frac{1}{2} \Big|.
\end{align*}
Regarding the necessary condition of this statement, in \cite{Peral1980} the author used elements of the form of $F_{\lambda}$ defined in \ref{def:g-lambda} (which periodization yields $f_{\lambda}$, see Section \ref{sec:localization}) as counterexample for $L^p$-based derivative regularity propagation. 
\hfill \break

Our goals are to extend this loss of $L^p$-based regularity to the corresponding periodic nonlinear flow with $d = 3$ and suitable $p \neq 2$ through a periodization of such elements $F_{\lambda}$; and from this to prove that $S_{\lambda}$ is a dense subset of $W^{\sigma-1,p}(\mathbb{T}^3)$ which does not preserve $W^{\sigma,p}$ regularity under the flow of \eqref{eq:cauchywave}, for $p > 2$ large enough. This is pathological in the sense that a proper randomization of initial data (see \eqref{eq:randomid}) ensures this propagation almost surely with respect to the Gaussian measure induced by such random data, as it was proved in \cite{GOTW22} by showing the quasi-invariance of this measure under the flow of \eqref{eq:cauchywave}. 


The (deterministic) analysis of these simple functions $F_{\lambda}$ let us identify a range of $p$, far enough from $2$, for which the linear wave evolution does not propagate the aforementioned $L^p$-based regularity on $\mathbb{R}^d$. Furthermore, making this analysis explicit allows to use a periodization argument in order to extend the same result to $\mathbb{T}^d$. Notice that, if $p = 2$, we recover the usual conservation of energy. A deeper analysis of this family of functions can be found in \cite{W65}. 

Once we identify the range of $p$ and $\lambda$ with loss of regularity for $S(t)(0,f_{\lambda})$ and suitable $t \neq 0$ (see Theorem \ref{thm:nonpropag-linear} below), in order to extend this behaviour to the three-dimensional nonlinear flow it suffices to bound the quantity ($T > 0$)
	\begin{align}\label{eq:quantity-to-control}
		\left\| \int_0^t \frac{\sin((t-\tau)|D|)}{|D|} u^3(\tau) d\tau \right\|_{L^{\infty}_T W^{\sigma,p}(\mathbb{T}^3)}.
	\end{align} 
	
	Namely, we look for an upper bound which depends essentially on the $H^r(\mathbb{T}^3)$ norm of the solution in a suitable way, for values of $r$ such that the initial datum is in $\vec{H}^r(\mathbb{T}^3) := H^r(\mathbb{T}^3) \times H^{r-1}(\mathbb{T}^3)$. This will let us apply the well-posedness theory for data on these spaces, see for instance \cite{Tzvetkov2019-notes}. 
	
All in all, this work deals with deterministic analysis, although the motivation (given in  Section \ref{sec:prob}) is deeply probabilistic. Our main contribution is the following.

\begin{theorem}\label{thm:nonpropag-nonlinear}
Let $p \geq 3$, $\lambda \in (3-\frac{3}{p},3)$ and $\sigma > 1$. Then $S_{\lambda}$ is a dense subset of $W^{\sigma-1,p}(\mathbb{T}^3)$ satisfying the following property:\\

For any $g \in S_{\lambda}$, the Cauchy problem \eqref{eq:cauchywave} admits a unique global solution $u$ in $C(\mathbb{R},H^{\sigma + \frac{1}{2}}(\mathbb{T}^3))$ such that, for any $t \in (-1,1) \setminus \{0\}$,
\begin{align*}
u(\cdot,t) \notin W^{\sigma,p}(\mathbb{T}^3).
\end{align*}
\end{theorem}
	
As we explained before, Theorem \ref{thm:nonpropag-nonlinear} will essentially follow from the next result, which is based on the analysis of the linear flow of \eqref{eq:cauchywave} on $\mathbb{R}^d$ and a periodization argument. Regarding the Euclidean linear case, we only make explicit the analysis from \cite{Peral1980} and \cite{Miyachi1980}.
	\begin{theorem}\label{thm:nonpropag-linear}
		Let $d \geq 2$, $p \neq 2$ and $\lambda \in (\max\{n-\frac{d}{p},\frac{d-1}{2}\},d)$ (satisfying $\frac{d-1}{2} - \lambda \notin - \mathbb{N}$) such that one of the following conditions holds:
\begin{itemize}
\item $p \geq \frac{2d}{d-1}$.
\item $p < \frac{2d}{d+1}$ and $\lambda \leq \frac{d+1}{2} - \frac{1}{p}$.  
\end{itemize}
Then $f_{\lambda} \in W^{\sigma-1,p}(\mathbb{T}^d)$ and
		\begin{align*}
			S(t)(0,f_{\lambda}) \notin W^{\sigma,p}(\mathbb{T}^d)
		\end{align*}
		for any $t \in (-1,1) \setminus \{0\}$ if $p \geq \frac{2d}{d-1}$, and for any $t \in (-1/2,1/2) \setminus \{0\}$ if $p < \frac{2d}{d+1}$.
	\end{theorem}


In Section \ref{sec:prob}, we provide the positive result after a suitable randomization of the initial data, using the work \cite{GOTW22}. In Section \ref{sec:linearcex}, following \cite{Peral1980}, we prove that the elements $F_{\lambda}$ for $x$ in $\mathbb{R}^d$ with $d > 1$ serve as counterexample for the propagation of regularity under the linear flow in $W^{\sigma,p}(\mathbb{R}^d)$ spaces. In Section \ref{sec:localization} we give a periodization argument to see the validity of this non-propagation for $f_{\lambda}$. In Section \ref{sec:proof-nonlinearcex} we extend this phenomenon to the nonlinear case with $d=3$ and $p \geq 3$ with the proof of Theorem \ref{thm:nonpropag-nonlinear}. Furthemore, in appendices \ref{app:bessel} and \ref{app:temp-dist-xalpha+} we provide some auxiliary results concerning Bessel functions and homogeneous distributions, which will be used throughout this work.

\begin{remark}
The bounds $1$ and $1/2$ for $|t|$ in Theorems \ref{thm:nonpropag-nonlinear} and \ref{thm:nonpropag-linear} are technical constraints due to the way we apply the finite speed of propagation from Proposition \ref{prop:finite-speed-prop}, and the method we use to show the divergence of
\begin{align*}
\| S(t)(0,f_{\lambda})\|_{W^{\sigma,p}(\mathbb{T}^d)},
\end{align*}
for $d \geq 2$ and suitable $\lambda$, $\sigma$ and $p$, after the periodization \eqref{eq:flambdadef-t}. Although they may be extended to any $t$, we admit this limitation since it does not prevent us from illustrating the aforementioned pathological behaviour: under $S(t)$ and, eventually, under the flow of \eqref{eq:cauchywave}, the solution ceases instantaneously to be in $W^{\sigma,p}(\mathbb{T}^3)$.
\end{remark}

\begin{remark}
This work shares the same philosophy than others in which the conceptual contrast between topological and measure-theoretic properties around the flow of certain PDEs is studied. We refer to \cite{ST20}, where they present the existence of a dense set of the Sobolev space of super-critical regularity as data which, convoluted with some approximate identity, generate a family of smooth solutions of \eqref{eq:cauchywave} that does not converge in the space of super-critical Sobolev regularity due to a norm-inflation mechanism; or in \cite{CG23}, where an analogous work is given for super-critical nonlinear Schrödinger equations.
%or in \cite{Ha14} \cite{GG15} \cite{GG17}, where they give examples of $G_{\delta}$ dense sets which give solutions of Hamiltonian PDE’s with growing Sobolev norms for large time,
\end{remark}

                                                                                                                                                                                                                                                                                                                                                                                                                                                                                                                                                                                                                                                                                                                                                                                                                                                                                                                                                                                                                                                                                                                                                                                                                                                                                                                                                                                                                                                                                   \section{Propagation of almost sure regularity under a Gaussian measure along the flow} \label{sec:prob}
We know the there is global well-posedness of the Cauchy problem \eqref{eq:cauchywave} if the initial datum belongs to $\vec{H}^r(\mathbb{T}^3)$, for $r \geq 1$ (see \cite{Tzvetkov2019-notes}), i.e. it admits a global flow $\Phi_{\text{NLW}}$ on these spaces. In order to be more precise with the notation of the flow, given $t \in \mathbb{R}$ we write
\begin{align*}
\Phi_{\text{NLW}}(t) :& \quad \vec{H}^r(\mathbb{T}^3) \rightarrow \vec{H}^r(\mathbb{T}^3) \\
&\quad (u_0 , v_0) \mapsto (u(t),\partial u(t)),
\end{align*}
where $(u,\partial_t u)$ is the global solution of the Cauchy problem \eqref{eq:cauchywave} with $\left. (u,\partial_t u)\right|_{t=0} = (u_0,v_0)$, $u_0 \in H^r(\mathbb{T}^3)$ and $u_1 \in H^{r-1}(\mathbb{T}^3)$, $r \geq 1$.

Given $s \in \mathbb{R}$, let $\vec{\mu}_s$ be the centered Gaussian measure with Cameron-Martin space $\vec{H}^s(\mathbb{T}^3)$. Denote $\vec{u} = (u,v)$. Such Gaussian measure has a formal density of the form

\begin{align*}
&d \vec{\mu}_s = Z_s^{-1} e^{-\frac{1}{2} \| \vec{u} \|^2_{\vec{H}^s}} d\vec{u} \nonumber \\
&= \prod_{k \in \mathbb{Z}^3} Z_{s,k}^{-1} e^{-\frac{1}{2} \langle k \rangle^{2s} |\widehat{u}(k)|^2} e^{-\frac{1}{2} \langle k \rangle^{2(s-1)} |\widehat{v}(k)|^2} d \widehat{u}(k) d\widehat{v}(k).
\end{align*}
Samples $\vec{u}^{\omega} = (u^{\omega},v^{\omega})$ from $\vec{\mu}_s$ can be constructed via the Karhunen-Loève expansions
\begin{align}\label{eq:randomid}
u^{\omega}(x) = \sum_{k \in \mathbb{Z}^3} \frac{g_k^{\omega}}{\langle k \rangle^s} e^{ikx} \text{ and }  v^{\omega}(x) = \sum_{k \in \mathbb{Z}^3} \frac{h_k^{\omega}}{\langle k \rangle^{s-1}} e^{ikx},
\end{align}
for $\{g_k^{\omega}\}_{k \in \mathbb{Z}^3}$ and $\{h_k^{\omega}\}_{k \in \mathbb{Z}^3}$ sequences of standard and independent gaussian random variables on $(\Omega,\mathbb{P})$ such that $g_k^{\omega} = \overline{g_{-k}^{\omega}}$ and $h_k^{\omega} = \overline{h_{-k}^{\omega}}$, in order to preserve values in $\mathbb{R}$. One can prove that these series converge in $L^2(\Omega;\vec{H}^r(\mathbb{T}^3))$ as long as
\begin{align*}
r < s - \frac{3}{2}
\end{align*}
and thus the map
\begin{align*}
\omega \in \Omega \mapsto (u^{\omega},v^{\omega})
\end{align*}
induces the Gaussian measure $\vec{\mu}_s$ as a probability measure with topological support $\vec{H}^r(\mathbb{T}^3)$ for the same range of $r$. Namely, both $u^{\omega}$ and $v^{\omega}$ converge $w$-a.s. in $H^r(\mathbb{T}^3)$ and $H^{r-1}(\mathbb{T}^3)$. We present now a result on quasi-invariance of this measure under the flow $\Phi_{\text{NLW}}(t)$ from \cite{GOTW22}.

\begin{theorem}\label{thm-qinv}
Let $s \geq 3$ be an odd integer. Then, $\vec{\mu}_s$ is quasi-invariant under the dynamics of the defocusing cubic $NLW$ from \eqref{eq:cauchywave}. Namely, for any $t \in \mathbb{R}$, the Gaussian measure $\vec{\mu}_s$ and its pushforward under $\Phi_{\text{\text{NLW}}}(t)$ are mutually absolutely continuous.
\end{theorem}

This ensures the propagation of almost sure properties with respect to the Gaussian measure along $\Phi_{\text{NLW}}$. In fact, in infinite dimensions this becomes important because many relevant properties regarding the small-scale behaviour under a Gaussian measure hold true with probability $1$ or $0$ (see \cite{Bogachev1998}). This is a consequence of Fernique's theorem (see for instance Theorem 2.7 in \cite{PZ14}): under a Gaussian measure, any given norm is finite with probability $1$ or $0$. 

Denote $\vec{W}^{\sigma,p}(\mathbb{T}^3) = W^{\sigma,p}(\mathbb{T}^3) \times W^{\sigma-1,p}(\mathbb{T}^3)$ for $\sigma \geq 1$ and $1 \leq p \leq \infty$. We prove that $\vec{\mu}_s(\vec{W}^{\sigma,p}(\mathbb{T}^3)) = 1$ if $\sigma < s - \frac{3}{2}$ at the end of the Section (see Proposition \ref{prop:asprop}), preceded by a useful result whose proof can be found in \cite{BurqTzvet}. By the quasi-invariance given in Theorem \ref{thm-qinv}
\begin{align*}
\vec{\mu}_s(\Phi_{\text{NLW}}(t)(\vec{W}^{\sigma,p}(\mathbb{T}^3))) = 1
\end{align*}
for any $t \in \mathbb{R}$, $\sigma < s - \frac{3}{2}$ and $1 \leq p \leq \infty$. 

Consider $\sigma > 1$ and $p\geq 3$, and let $s > \max\{3,\sigma + \frac{3}{2}\}$. Thus, the regularity in terms of these $\vec{W}^{\sigma,p}(\mathbb{T}^3)$ spaces is $\vec{\mu}_s$-almost surely propagated, in contrast with Theorem \ref{thm:nonpropag-nonlinear}, where we proved that, for any $\lambda \in (3 - \frac{3}{p},3)$, the pathological set $S_{\lambda} \subset W^{\sigma-1,p}(\mathbb{T}^3)$ of initial velocities in \eqref{eq:cauchywave} is dense in $W^{\sigma-1,p}(\mathbb{T}^3)$ and gives 
\begin{align*}
\Phi_{\text{NLW}}(t)(\{0\} \times S_{\lambda}) \cap \vec{W}^{\sigma,p}(\mathbb{T}^3) = \emptyset
\end{align*}
for any $t \in (-1,1) \setminus \{0\}$.
\begin{lemma}\label{lemma:1deg-wchaos}
		Let $d \geq 1$, $(a_k) \in \l^2(\mathbb{Z}^d)$ and $(g^{\omega}_k)$ a sequence of complex, pairwise independent, standard Gaussian random variables . Then, the random variable
		\begin{align*}
			w \mapsto \sum_{k \in \mathbb{Z}^d} a_k g^{\omega}_k
		\end{align*}
		is sub-Gaussian (up to normalization to get unit variance). In particular, there exists $C>0$ such that, for every $r \geq 1$,
		\begin{align}
			\| \sum_{k \in \mathbb{Z}^d} a_k g^{\omega}_k\|_{L^r_w} \leq C \sqrt{r} \|(a_k)\|_{\l^2(\mathbb{Z}^d)}.
		\end{align}
		\label{lemma:randomsum-key}
	\end{lemma}

\begin{proposition}\label{prop:asprop}
Let $s \in \mathbb{R}$, $1 \leq p \leq \infty$ and $\sigma < s - \frac{3}{2}$. Then $(u^{\omega},v^{\omega}) \in \vec{W}^{\sigma,p}(\mathbb{T}^3)$, $\omega$-almost surely. In other words, $\vec{\mu}(\vec{W}^{\sigma,p}(\mathbb{T}^3)) = 1$.
\end{proposition}
\begin{proof}
Let $p \in [1,\infty)$. If $p \leq Q < \infty$, by Minkowski inequality and Lemma \ref{lemma:1deg-wchaos}
\begin{align*}
\| \| u^{\omega} \|_{W^{\sigma,p}(\mathbb{T}^3)} \|_{L^Q_w} \lesssim \| \| \langle D \rangle^{\sigma} u^{\omega} \|_{L^Q_{\omega}} \|_{L^p(\mathbb{T}^3)} \lesssim \sqrt{Q} \| (\langle k \rangle^{\sigma - s})_{k \in \mathbb{Z}^3} \|_{\l^2(\mathbb{Z}^3)},
\end{align*}
which converges if $\sigma < s - \frac{3}{2}$. If $Q < p$, by Hölder's inequality
\begin{align*}
\| \| u^{\omega} \|_{W^{\sigma,p}(\mathbb{T}^3)} \|_{L^Q_{\omega}} \lesssim \| \| \langle D \rangle^{\sigma} u^{\omega} \|_{L^p_{\omega}} \|_{L^p(\mathbb{T}^3)} \lesssim \sqrt{p} \| (\langle k \rangle^{\sigma - s})_{k} \|_{\l^2(\mathbb{Z}^3)},
\end{align*}
which again converges if $\sigma < s - \frac{3}{2}$. Thus, for any $Q \in [1,\infty)$ and $\sigma < s - \frac{3}{2}$
\begin{align*}
\| \| u^{\omega} \|_{W^{\sigma,p}(\mathbb{T}^3)} \|_{L^Q_{\omega}} \lesssim \sqrt{Q} \left(\max_{Q \in [1,p]} \sqrt{p/Q}\right) = \sqrt{Qp}.
\end{align*}
By Markov inequality and proceeding in an identical way with $v^{\omega}$, this implies that there exists some suitable constants $C,c>0$ independent of $\lambda$ such that
\begin{align*}
\vec{\mu}_s(\|(u^{\omega},v^{\omega})\|_{\vec{W}^{\sigma,p}} > \lambda) \leq C e^{-c\lambda^2/p} \text{ for any } \lambda > 0.
\end{align*}
In other words, given $\varepsilon \in (0,1)$
\begin{align*}
\| (u^{\omega},v^{\omega}) \|_{\vec{W}^{\sigma,p}(\mathbb{T}^3)} \lesssim (-p\log(\varepsilon))^{1/2}
\end{align*}
with probability $1 - \varepsilon$. In order to pass to probability $1$, denote
\begin{align*}
A_N = \{\omega \in \Omega : \| (u^{\omega},v^{\omega}) \|_{\vec{W}^{\sigma,p}(\mathbb{T}^3)}> N\}, \quad N \in \mathbb{N}.
\end{align*}
We obtain that
\begin{align*}
\sum_{N \in \mathbb{N}} \mathbb{P}(A_N) \lesssim \sum_{N \in \mathbb{N}} e^{-cN^2/p} < \infty.
\end{align*}
By Borel-Cantelli Lemma, this implies that 
\begin{align*}
\mathbb{P}(\limsup_{N \rightarrow \infty} A_N) = 0 \Longrightarrow \mathbb{P}(\liminf_{N \rightarrow \infty} A_N^c) = 1,
\end{align*}
i.e. for almost every $\omega \in \Omega$ there exists some $N_{\omega} \in \mathbb{N}$ such that
\begin{align*}
\| (u^{\omega},v^{\omega}) \|_{\vec{W}^{\sigma,p}(\mathbb{T}^3)} \leq N_{\omega}.
\end{align*}
This concludes the proof for the case $p < \infty$. 

Regarding $p=\infty$, with the operator $D = \frac{1}{i}\nabla$ adapted to functions on $\mathbb{T}^3$, given $j \geq 1$ consider the projections $Q_j = \varphi_j(D)$, $Q_0 = \varphi_0(D)$, with $\varphi_0, \varphi_j \in C^{\infty}(\mathbb{R}^3)$ and $\supp(\varphi_0) \subset \overline{B}(0,1)$, $\supp(\varphi_j) \subset \overline{B}(0,2^{j+1})\setminus \overline{B}(0,2^{j-1})$. Since $\sigma < s - \frac{3}{2}$, we know that there exists some $q \in [1,\infty)$ such that $\sigma < s - \frac{3}{2} - \frac{1}{q}$. Consider $Q \geq 1$, so that by Bernstein inequality and a similar procedure to the case $p < \infty$,
\begin{align*}
&\| \| Q_j u^{\omega} \|_{W^{\sigma,\infty}} \|_{L^Q_{\omega}} \lesssim 2^{\frac{j}{q}} \| \| Q_j u^{\omega} \|_{W^{\sigma,q}} \|_{L^Q_{\omega}} \lesssim \sqrt{Qq} 2^{j\left( \frac{1}{q} + \sigma - s + \frac{3}{2} \right)}, \nonumber \\
&\| \| Q_0 u^{\omega} \|_{W^{\sigma,\infty}} \|_{L^Q_{\omega}} \lesssim \sqrt{Qq}.
\end{align*}
Let $(\sigma_j)_{j \in \mathbb{N} \cup \{0\}} \subset (0,\infty)$ be a sequence in the unit ball of $\l^1(\mathbb{N})$ such that\footnote{It would suffice to take, for each $j \in \mathbb{N} \cup \{0\}$,
\begin{align*}
\sigma_j := c_{\sigma}^{-1} 2^{-\frac{j}{2}(s -  \sigma - \frac{3}{2} - \frac{1}{q})}, \quad c_{\sigma} = \sum_{j \in \mathbb{N} \cup \{0\}} 2^{-\frac{j}{2}(s - \sigma - \frac{3}{2} - \frac{1}{q})}.
\end{align*}}
\begin{align*}
\sigma_j \geq 2^{-j(s -  \sigma - \frac{3}{2} - \frac{1}{q})}, \quad j \geq j_0,
\end{align*}
for some $j_0 \in \mathbb{N} \cup \{0\}$. This choice of $\sigma_j$ let us take $C,c > 0$ independent of $j$ such that (doing also the analogous computation for $v^{\omega}$)
\begin{align*}
&\vec{\mu}_s(\| (Q_j u^{\omega},Q_jv^{\omega}) \|_{\vec{W}^{\sigma,\infty}(\mathbb{T}^3)} > \sigma_j \lambda) \leq C e^{-\frac{c\lambda^2 \sigma_j^2}{q} 2^{j(2s - 2 \sigma - 3 - \frac{2}{q})}}, \nonumber \\
&\vec{\mu}_s(\| (Q_0 u^{\omega},Q_0 v^{\omega}) \|_{\vec{W}^{\sigma,\infty}(\mathbb{T}^3)} > \sigma_0 \lambda) \leq C e^{-\frac{c\lambda^2\sigma_0^2}{q}}
\end{align*}
Summing all the pieces,
\begin{align*}
&\vec{\mu}_s(\| (u^{\omega},v^{\omega}) \|_{\vec{W}^{\sigma,\infty}} > \lambda) \nonumber \\ 
&\leq \vec{\mu}_s(\| (Q_0 u^{\omega},Q_0v^{\omega}) \|_{\vec{W}^{\sigma,\infty}} > \sigma_0 \lambda) + \sum_{j \in \mathbb{N}} \mu_s(\| (Q_j u^{\omega},Q_jv^{\omega}) \|_{\vec{W}^{\sigma,\infty}} > \sigma_j \lambda) \\
&\lesssim \sum_{j \in \mathbb{N}_0} e^{-\frac{c\lambda^2\sigma_j^2}{q} 2^{j(2s - 2 \sigma - 3 - \frac{2}{q})}} \lesssim e^{-c_1\lambda^2/q}
\end{align*}
for some constant $c_1 > 0$ independent of $\lambda$. Therefore, using Borel-Cantelli Lemma in the same way as above, we have that 
\begin{align*}
(u^{\omega},v^{\omega}) \in \vec{W}^{\sigma,\infty}(\mathbb{T}^3), \quad \omega\text{-almost surely}.
\end{align*}
\end{proof}


\section{Deterministic counterexample for $S(t)$ on $\mathbb{R}^d$}\label{sec:linearcex}	


Consider the space variable as $x \in \mathbb{R}^d$. Let $\sigma \geq 1$ and  $\mathcal{F} : \mathcal{S}'(\mathbb{R}^d) \rightarrow \mathcal{S}'(\mathbb{R}^d)$ be the distributional Fourier transform map. Let $\psi$ be a symmetric mollifier, i.e. an element of $C^{\infty}_c(\mathbb{R}^d)$ such that
\begin{align*}
&\bullet \int_{\mathbb{R}^d}\psi(x)dx = 1, \\
&\bullet \lim_{\varepsilon \rightarrow 0^+} \varepsilon^{-d} \psi(\cdot/\varepsilon) = \delta \text{ in } \mathcal{S}'(\mathbb{R}^d), \\
&\bullet \psi(\xi) = \psi_1(|\xi|) \text{ for any } \xi \in \mathbb{R}^d,
\end{align*}
where $\psi_1 : [0,\infty) \rightarrow \mathbb{R}$ is smooth (chosen to have support contained in $[0,1]$) and $\delta$ is the Dirac delta distribution. Denote, for any $\xi \in \mathbb{R}^d$, $\psi_{1/4}(\xi) = 4^d \psi(4 \xi)$ and
\begin{align}\label{eq:def-h}
h(\xi) = 1 - (\psi_{1/4} \ast \mathrm1_{B(0,3/4)})(\xi),
\end{align}
where $\mathrm1_{B(0,3/4)}$ is the characteristic function of the ball centered at $0$ with radius $3/4$ in $\mathbb{R}^d$. Therefore, $h$ is a smooth function on $\mathbb{R}^d$ which is $0$ on $B(0,1/2)$ (as well as all its derivatives), and which is $1$ on $B(0,1)^c$. This setting allows us to define the main objects of this Section.

\begin{definition}\label{def:g-lambda}
Let $d \in \mathbb{N}$ and $\lambda \in (0,d)$. Define $F_{\lambda} = \mathcal{F}^{-1}[g_{\lambda}]$, where $g_{\lambda}(\xi) = \langle \xi \rangle^{-\sigma+1} h(\xi)|\xi|^{-\lambda}$.
\end{definition}
The main results of this Section are the following ones. Note that $F_{\lambda} \in \mathcal{S}'(\mathbb{R}^d)$ for any $\lambda \in \mathbb{R}$.

\begin{proposition}\label{prop:Flambda-inWsigmap}
Let $p\in[1,\infty)$. The elements $F_{\lambda}$ given in Definition \ref{def:g-lambda} belong to $W^{\sigma-1,p}(\mathbb{R}^d)$ if $\lambda \in (d - \frac{d}{p},d)$.
\end{proposition}


\begin{theorem}\label{thm:nonpropag-linear-rn}
Let $d \geq 2$, $1 \leq p < \infty$, $t \neq 0$ and $\lambda \in (\max\{ \frac{d-1}{2} , d - \frac{d}{p}\},d)$ satisfying $\frac{d-1}{2} - \lambda \notin -\mathbb{N}$. Then $F_{\lambda} \in W^{\sigma-1,p}(\mathbb{R}^d)$ and $\frac{\sin(t|D|)}{|D|} F_{\lambda} \notin W^{\sigma,p}(\mathbb{R}^d)$ if one of the following conditions holds:
\begin{itemize}
\item $p \geq \frac{2d}{d-1}$.
\item $p < \frac{2d}{d+1}$ and $\lambda \leq \frac{d+1}{2} - \frac{1}{p}$.  
\end{itemize}
\end{theorem}


Proposition \ref{prop:Flambda-inWsigmap} will be a straightforward consequence of Lemma \ref{lemma:Flambda} below, where we characterize the functions $F_{\lambda}$. In order to prove Theorem \ref{thm:nonpropag-linear-rn}, we provide some preliminary results that will be useful later to analyze $\langle D \rangle^{\sigma} \frac{e^{it|D|}}{|D|} F_{\lambda}$. Namely, Lemma \ref{lemma:ftransformradial_bessel} below will let us use an integral representation of radial elements in $L^p(\mathbb{R}^d)$, $1 \leq p < \infty$, in terms of Bessel functions, and Lemma \ref{lemma:invft-xalpha-nu} gives us the explicit form for the distributional inverse Fourier transforms of certain homogeneous distributions that will come later (its proof will follow the strategy from \cite{GS64}). For the sake of completeness, we include their proofs in Appendices \ref{app:bessel} and \ref{app:temp-dist-xalpha+}. 

\begin{remark}
For both cases in $\mathbb{R}^d$ and $\mathbb{T}^d$, the main computation to understand whether $S(t)(0,F_{\lambda})$ (resp. $S(t)(0,f_{\lambda})$) belongs to $W^{\sigma,p}(\mathbb{R}^d)$ (resp. $W^{\sigma,p}(\mathbb{T}^d)$) lies in the operator $\langle D \rangle^{\sigma} \frac{e^{\pm it|D|}}{|D|}$. From now on, we will mainly consider the term corresponding to $e^{it|D|}$ for simplicity, since the analysis with $e^{-it|D|}$ is analogous.
\end{remark}


\begin{lemma}\label{lemma:Flambda}
Let $d \in \mathbb{N}$, $x \in \mathbb{R}^d \setminus \{0\}$ and $\lambda \in (0,d)$. Then
\begin{align*}
\langle D \rangle^{\sigma-1} F_{\lambda}(x) = \gamma(d,\lambda) |x|^{-n+\lambda} + G_{\lambda,d}(x),
\end{align*}
where $G_{\lambda,d} \in C^{\infty}(\mathbb{R}^d)$ and $\gamma(d,\lambda)$ is a constant only depending on $n$ and $\lambda$. On the other hand,
\begin{align*}
|\langle D \rangle^{\sigma-1} F_{\lambda}(x)| \lesssim_{d,\lambda} |x|^{-N}
\end{align*} 
for any $N > d - \lambda$.
\end{lemma}
\begin{proof}
Let $M > 0$ and $x \in B(0,M)^c$. Let $\varphi \in C^{\infty}_c(\mathbb{R}^d\setminus \{0\})$ with $\supp \varphi \subset \{\xi \in \mathbb{R}^d : \frac{1}{2} \leq |\xi| \leq 2\}$ satisfying
\begin{align*}
\sum_{j = -\infty}^{\infty} \varphi(2^{-j}\xi) = 1, \quad \xi \neq 0,
\end{align*}
so that for any $\xi \neq 0$ there are at most two terms in this sum which are nonzero (see page $197$ of \cite{MS1} for details). This allows us to consider the Littlewood-Paley decomposition
\begin{align}
			F_{\lambda} = \sum_{j \geq -1} P_{j}F_{\lambda},
			\label{eq:lp-decomp}
		\end{align}
		where $P_j = \varphi(D/2^j)$, and where we used that $g_{\lambda}(\xi) = 0$ if $|\xi| \leq \frac{1}{2}$. Denote $B_j = \overline{B}(0,2^{j+1}) \setminus \overline{B}(0,2^{j-1})$ for $j \in \mathbb{Z}$. Then
\begin{align*}
|\langle D \rangle^{\sigma-1} F_{\lambda}(x)| \leq \sum_{j \geq -1} \left| \int_{B_j} \varphi(\xi/2^j) g_{\lambda}(\xi) e^{i\xi x} d\xi \right| = \sum_{j\geq -1} |I_j|.
\end{align*}
By doing a change of variables $\xi = 2^j \eta$, we obtain
\begin{align}\label{eq:oscint-0}
&I_j = 2^{j(d - \lambda)} \int_{B_0} h(2^j \eta) \varphi(\eta) |\eta|^{-\lambda} e^{i(2^j |x|) \eta \cdot \frac{x}{|x|}} d\eta, \quad j \geq -1.
\end{align}
Let $j \geq 1$. Then \eqref{eq:oscint-0} is an oscillatory integral with a phase $(2^j |x|)\phi(\eta)$ for which $\nabla_{\eta} \phi(\eta) = x/|x|$, which is not zero because $|x| \geq M$ and which satisfies $|\nabla_{\eta}\phi(\eta)| = 1$ for any $\eta \in B_0$. Namely
\begin{align}\label{eq:oscint-o1}
|I_j| = 2^{j(d - \lambda)} \left| \int_{B_0} \varphi(\eta) |\eta|^{-\lambda} e^{i(2^j |x|) \eta \cdot \frac{x}{|x|}} d\eta \right|,
\end{align}
where $\varphi(\eta)|\eta|^{-\lambda}$ is smooth and compactly supported in $B_0$. Then, for any $N \in \mathbb{N}$,
\begin{align*}
|I_j| \lesssim_{d,\lambda,\varphi} |x|^{-N} 2^{j(d - \lambda - N)}.
\end{align*}
Regarding $j \in \{-1,0\}$, note that $\varphi(\eta) |\eta|^{-\lambda} h(2^j \eta)$ is also smooth and compactly supported in $B_0$, so we can obtain \eqref{eq:oscint-o1} too. All in all,
\begin{align}\label{eq:nonstatphase-oscint}
|\langle D \rangle^{\sigma-1} F_{\lambda}(x)| \lesssim |x|^{-N} \left( \sum_{j \geq -1} 2^{j(d-\lambda-N)} \right),
\end{align}
which requires $N > d - \lambda$ to be convergent at the right hand side. \\

On the other hand, we have
\begin{align*}
g_{\lambda}(\xi) = \langle \xi \rangle^{-\sigma+1}|\xi|^{-\lambda} +  \langle \xi \rangle^{-\sigma+1} |\xi|^{-\lambda} ( h(\xi) - 1 )
\end{align*}
in $\mathcal{S}'(\mathbb{R}^d)$ (recall that $\lambda \in (0,d)$), so
\begin{align}\label{eq:Flambda-proof1}
\langle D \rangle^{\sigma-1} F_{\lambda} = \mathcal{F}[|\xi|^{-\lambda}] + G_{\lambda,d}
\end{align}
where
\begin{align*}
G_{\lambda,d}(x) = \int_{B(0,1)} e^{i\xi x} |\xi|^{-\lambda}(h(\xi) - 1)d\xi.
\end{align*}
Then, by the dominated convergence theorem, $G_{\lambda,d}$ is infinitely differentiable as long as $\lambda < n$. Regarding the first term, from standard theory\footnote{In fact, if we take Gaussian test functions and perform standard computations, we obtain
\begin{align*}
\gamma(d,\lambda) = \pi^{\lambda-\frac{d}{2}}\frac{\Gamma(\frac{d-\lambda}{2})}{\Gamma(\frac{\lambda}{2})}
\end{align*}} on homogeneous distributions one can prove that $\mathcal{F}[|\xi|^{-\lambda}] = \gamma(d,\lambda) |x|^{\lambda-n}$ in $\mathcal{S}'(\mathbb{R}^d)$ (see \cite{wolff} for  details), where $\gamma(d,\lambda)$ is some constant.
Therefore
\begin{align*}
\langle D \rangle^{\sigma-1} F_{\lambda}(x) = \gamma(d,\lambda) |x|^{-n+\lambda} + G_{\lambda,d}(x)
\end{align*}
in $\mathcal{S}'(\mathbb{R}^d)$. The modulus of this last expression is less than $|x|^{-N}$, $N > d - \lambda$, for $|x| < M$ with $M > 0$ small enough. This concludes the proof.
\end{proof}

\begin{lemma}\label{lemma:ftransformradial_bessel}
		Let $p \in [1,\infty)$. Suppose $f$ is a radial function in  $L^p(\mathbb{R}^d)$, $d \geq 2$; thus, $f(x) = f_0(|x|)$ for a.e. $x \in \mathbb{R}^d$. Moreover, assume that the integral
\begin{align}\label{eq:convcond-radialexp}
\int_0^{\infty} f_0(\rho) J_{\frac{d-2}{2}}(\rho|\xi|) \rho^{d/2} d\rho
\end{align}		
converges in a.e. $\xi \in \mathbb{R}^d$. Then the Fourier transform $\mathcal{F}[f]$ is given also by a radial function, and it has the form
		\begin{align*}
			\mathcal{F}[f](\xi) = c(d) |\xi|^{-\frac{d-2}{2}} \int_0^{\infty} f_0(\rho) J_{\frac{d-2}{2}}(\rho|\xi|) \rho^{d/2} d\rho \quad \text{for a.e. } \xi \in \mathbb{R}^d,
		\end{align*}
		where $c(d)$ is some constant depending only on the dimension.
	\end{lemma}
	\begin{remark}
	Lemma \ref{lemma:ftransformradial_bessel} will be applied to $f = e^{it|\xi|} \langle \xi \rangle^{\sigma-1} g_{\lambda}(\xi)$ in the proof of Theorem \ref{thm:flambda-tnotzero}. Hypothesis \eqref{eq:convcond-radialexp} will be implicitly given in such a proof.
	\end{remark}
	\begin{proof}
	See Appendix \ref{app:bessel}.
	\end{proof}

We denote $x_+$ and $x_-$ the functions $x \mathbb{I}_{(0,\infty)}(x)$ and $-x \mathbb{I}_{(-\infty,0)}(x)$ defined on $\mathbb{R}$. The proof of the following Lemma is given in Appendix \ref{app:temp-dist-xalpha+}.
	
	\begin{lemma}\label{lemma:invft-xalpha-nu}
	Let $\alpha \in \mathbb{C}$ such that $\text{Re}(\alpha) \notin -\mathbb{N}$. Then, the tempered distributions $x^{\alpha} \mathbb{I}_{(1,\infty)}(x)$ admit the inverse distributional Fourier transform
	\begin{align*}
	&\mathcal{F}^{-1}[x^{\alpha} \mathbb{I}_{(1,\infty)}(x)](\theta)= J(\alpha) \{ e^{i\frac{\pi}{2}(\alpha+1)} \theta_+^{-\alpha-1} + e^{-i\frac{\pi}{2}(\alpha+1)} \theta_-^{-\alpha-1} \}
	\end{align*}
	where $J : \mathbb{C} \rightarrow \mathbb{C}$ is a function such that
	\begin{align*}
	|J(\alpha)| \leq \Gamma(\text{Re}(\alpha)+1) + \frac{\pi}{2}.
	\end{align*}
	\end{lemma}
	%\begin{corollary}\label{cor:invft-xalpha-nu}
%Let $\alpha \in \mathbb{C}$ such that $\text{Re}(\alpha) \neq -1,-2,-3,\dots$. Then, the elements $x_+^{\alpha} \in \mathcal{S}'(\mathbb{R})$ admit the following distributional Fourier transform:
%	\begin{align*}
%	\mathcal{F}[x_+^{\alpha}](\theta) = i \Gamma(\alpha+1) \{ e^{i\frac{\pi}{2}\alpha}\theta_+^{-\alpha-1} - e^{-i\frac{\pi}{2}\alpha} \theta_-^{-\alpha-1} \}. 
%	\end{align*}
%	Moreover, for any $\tau > 0$,
%	\begin{align*}
%	\mathcal{F}[x_+^{\alpha}e^{-\tau x}](\theta) = i \Gamma(\alpha+1) e^{i\frac{\pi}{2}\alpha} (\theta + i \tau)^{-\alpha-1}.
%	\end{align*}
%	\end{corollary}

Furthermore, we provide a result that will simplify the computation. Namely, that the non-convergence of 
\begin{align*}
\left\|\langle D \rangle^{\sigma} \frac{e^{it|D|}}{|D|} F_{\lambda} \right\|_{L^p(\mathbb{R}^d)}
\end{align*} 
follows from the non-convergence of 
\begin{align*}
\left\|\langle D \rangle^{\sigma-1} e^{it|D|} F_{\lambda} \right\|_{L^p(\mathbb{R}^d)}.
\end{align*}

\begin{lemma}\label{lemma:bdd-finitemeasure}
Let $1 \leq q \leq \infty$ and $F \in L^q(\mathbb{R}^d)$. Then
\begin{align*}
\left\| \frac{|D|}{\langle D \rangle} F \right\|_{L^q(\mathbb{R}^d)} \lesssim \|F\|_{L^q(\mathbb{R}^d)}.
\end{align*}
\end{lemma}
\begin{proof}
Note that $\frac{|D|}{\langle D \rangle}$ is an operator given by the Fourier multiplier $\frac{|\xi|}{(1 + |\xi|^2)^{1/2}}$. The latter is the Fourier transform of a finite measure $\beta$ in $\mathbb{R}^d$, see \cite{S70} (Lemma 2 in page 133) for the details. This implies that
\begin{align*}
\left\| \frac{|D|}{\langle D \rangle} F \right\|_{L^q(\mathbb{R}^d)} = \left\| \beta \ast F \right\|_{L^q(\mathbb{R}^d)} \lesssim_{\beta} \|F\|_{L^q(\mathbb{R}^d)}.
\end{align*}
\end{proof}

\begin{comment}
\begin{remark}
Assume that we have proved that $\|F\|_{L^q(\mathbb{R}^d)} = \infty$ for some $q \geq 1$. Thanks to result \eqref{lemma:bdd-finitemeasure}, we have that
\begin{align*}
\left\|\frac{\langle D \rangle}{|D|} F\right\|_{L^q(\mathbb{R}^d)} \gtrsim \| F \|_{L^q(\mathbb{R}^d)},
\end{align*}
so $\|\frac{\langle D \rangle}{|D|} F\|_{L^q(\mathbb{R}^d)} = \infty$ would follow from the first assumption. This let us simplify the analysis to $\langle D \rangle^{\sigma-1} e^{it|D|} F_{\lambda}$ in what follows.
\end{remark}
\end{comment}

The following result provides us with the explicit form of $\langle D \rangle^{\sigma-1} e^{it|D|} F_{\lambda}$, which will let us identify the conditions on $\lambda$ and $p$ for the non-convergence of $\|\frac{\sin(t|D|)}{|D|} F_{\lambda}\|_{W^{\sigma,p}(\mathbb{R}^d)}$ for any $t \neq 0$.
\begin{theorem}\label{thm:flambda-tnotzero}
Let $d>1$, $t \in \mathbb{R} \setminus \{0\}$, $\lambda \in (\frac{d-1}{2},d)$ such that $\frac{d-1}{2} - \lambda \notin -\mathbb{N}$, and $J$ the function given in Lemma \ref{lemma:invft-xalpha-nu}. Then the following holds:
\begin{itemize}
\item For a.e. $x \in \mathbb{R}^d \setminus \{0\}$ such that $|x| \neq t$ and $-|x| \neq t$:
\begin{align*}
\langle &D \rangle^{\sigma-1} e^{it|D|} F_{\lambda}(x) \nonumber \\
&= J\left(\frac{d-1}{2}-\lambda\right)  |x|^{-\frac{d-1}{2}} \left\{ c_+ [ e^{i\frac{\pi}{2}(\frac{d+1}{2} - \lambda)} (|x| + t)_+^{\lambda-\frac{d+1}{2}}\right. \\
&\hspace{48mm} + e^{-i\frac{\pi}{2}(\frac{d+1}{2} - \lambda)} \left(|x| + t\right)_-^{\lambda-\frac{d+1}{2}}] \nonumber \\
&\left. + c_-  [ e^{i\frac{\pi}{2}(\frac{d+1}{2} - \lambda)}(-|x| + t)_+^{\lambda-\frac{d+1}{2}} + e^{-i\frac{\pi}{2}(\frac{d+1}{2} - \lambda)} \left( -|x| + t\right)_-^{\lambda-\frac{d+1}{2}} ] \right\} \nonumber  \\
&+ c_{\lambda}(x) + D_{\lambda}(x),
\end{align*}
for $c_+$ and $c_-$ constants, $D_{\lambda} \in C^{\infty}(\mathbb{R}^d)$ and 
\begin{align*}
|c_{\lambda}(x)| \lesssim \min\{|x|^{-\frac{d+1}{2}} , |x|^{-\frac{d-1}{2}}\}
\end{align*}
if $\lambda \in (\frac{d+1}{2},d)$, and 
\begin{align*}
|c_{\lambda}(x)| \lesssim |x|^{-\frac{d+1}{2}}
\end{align*}
if $\lambda \in (\frac{d-1}{2},d)$.
%; and if $x = 0$,
%\begin{align*}
%&\langle D \rangle^{\sigma} e^{\pm i t |D|} F_{\lambda}(0) = w_{d-1}\{ \Gamma(d - \lambda) + J(n-\lambda-1,\nu) \} \nonumber \\
%& \times \{ c_+ e^{i\frac{\pi}{2}(n-\lambda)} t_+^{-n + \lambda} + c_- e^{-i\frac{\pi}{2}(n-\lambda)} t_-^{-n + \lambda}\},
%\end{align*}
%for $w_{d-1}$ the surface of the $(d-1)$-dimensional sphere.

\item  Moreover, for any $A > 1$, $x \in \mathbb{R}^d \setminus B(0,A|t|)$ and $N > d - \lambda$, we have that
\begin{align*}
|\langle D \rangle^{\sigma-1} e^{it|D|} F_{\lambda}(x)| \lesssim_{d,\lambda} |x|^{-N}.
\end{align*}

\end{itemize}
\end{theorem}

\begin{proof}
Let $A > 1$ and $x \in \overline{B}(0,A|t|)^c$. We can repeat the procedure from the Littlewood-Paley decomposition in \eqref{eq:lp-decomp} with a slight modification. Indeed, maintaining the same notation,
\begin{align*}
&|\langle D \rangle^{\sigma-1} e^{it|D|} F_{\lambda}(x)| \leq \sum_{j \geq -1} \left| \int_{B_j} e^{i \xi x + t |\xi|}(\varphi(\xi/2^j) |\xi|^{-\lambda} h(\xi)) d\xi \right| \nonumber \\
&= \sum_{j \geq 1} 2^{j(d-\lambda)} \left| \int_{B_0} e^{i 2^j |x| (\frac{\eta x + t |\eta|}{|x|})}(\varphi(\eta) |\eta|^{-\lambda} ) d\eta \right| \nonumber \\
&+ \sum_{j=-1}^0 2^{j(d-\lambda)} \left| \int_{B_0} e^{i 2^j |x| (\frac{\eta x + t |\eta|}{|x|})}(\varphi(\eta) |\eta|^{-\lambda} h(2^j \eta)) d\eta \right|
\end{align*}
The phase $\phi(\xi) = \frac{\eta x + t |\eta|}{|x|}$ does not have critical points for $|x| > |t|$. Indeed,
\begin{align*}
\nabla_{\eta} \phi(\eta) = \frac{1}{|x|} \left( x + \frac{t \eta}{|\eta|} \right),
\end{align*}
so $\nabla_{\eta} \phi(\eta) = 0$ would imply $|x| = |t|$. Moreover, since $|x| > A |t|$, we have that for any $\eta \in B_0$
\begin{align*}
|\nabla_{\eta} \phi(\eta)| \geq 1 - \frac{|t|}{|x|} > 1 - \frac{1}{A}.
\end{align*}
Moreover, $\varphi(\eta) |\eta|^{-\lambda} h(2^j \eta)$ is smooth and compactly supported in $B_0$ for any $j \geq -1$ (notice that $h(2^j \cdot) = 1$ for $j \geq 1$).  Thus we can do the standard integration by parts to obtain
\begin{align*}
|\langle D \rangle^{\sigma-1} e^{it|D|} F_{\lambda}(x)| \lesssim_{d,\lambda,\varphi} |x|^{-N} \left( \sum_{j \geq -1} 2^{j(d-\lambda-N)} \right),
\end{align*}
where $N > d - \lambda$ suffices for the convergence of the right-hand side. %Note that, in this procedure, we did not use the presence of $e^{-\gamma |\xi|}$ in the integrand.

Let $x \in \mathbb{R}^d \setminus \{0\}$ such that $|x| \neq t$ and $-|x| \neq t$. Since $\lambda < n$, there exists some $q > 1$ such that $\lambda > \frac{d}{q}$. Recalling Definition \ref{def:g-lambda}, this implies that 
\begin{align*}
\langle \xi \rangle^{\sigma-1} e^{it|\xi|} g_{\lambda}(\xi) \in L^q(\mathbb{R}^d).
\end{align*}
Then, for almost every $x \in \mathbb{R}^d$, from Lemma \ref{lemma:ftransformradial_bessel} we have that
\begin{align*}
&\langle D \rangle^{\sigma-1} e^{it|D|} F_{\lambda}(x) \nonumber \\
&=c(d) |x|^{-\frac{d-2}{2}} \int_1^{\infty}  e^{it\rho}  \rho^{-\lambda} J_{\frac{d-2}{2}}(\rho|x|) \rho^{d/2} d\rho + \int_{|\xi| \leq 1} e^{ix\xi + it|\xi|} |\xi|^{-\lambda} h(\xi) d\xi \nonumber \\
&=(i) + (ii).  
\end{align*}
The object $(ii)$ is in $C^{\infty}(\mathbb{R}^d)$ for each $t \neq 0$ and $\lambda < n$. We focus on $(i)$. From Bessel functions asymptotic expansion \eqref{eq:asympBessel}, we have that
\begin{align*}
J_{\frac{d-2}{2}}(s)  - \left( c_+ s^{-1/2} e^{is} + c_- s^{-1/2} e^{-is} \right) = R_{\frac{d-2}{2}}(s)
\end{align*}
for $R_{\frac{d-2}{2}}(s) = \mathcal{O}(s^{-3/2})$ when $s \rightarrow \infty$, and 
$c_+, c_- \in \mathbb{C}$ constants. This expansion and the Poisson representation of Bessel functions given in Lemma \ref{lemma:bessel_kreal} implies that, for any $s \in \mathbb{R} \setminus \{0\}$,
\begin{align}\label{eq:ourbessel-asymp}
|s^{\frac{1}{2}}R_{\frac{d-2}{2}}(s)| \lesssim \min\{ 1 , s^{-1} \}.
\end{align}
Observe that $\frac{d-2}{2} > -\frac{1}{2}$ if and only if $d > 1$, i.e. $d>1$ is required for a well definition of the Bessel functions as we have defined them. We write
\begin{align*}
&(i) = |x|^{-\frac{d-2}{2}} \left( c_+ \int_1^{\infty}  e^{i\rho|x| + it\rho} \rho^{-\lambda} (\rho|x|)^{-1/2} \rho^{d/2} d\rho \right.   \nonumber \\
&\left. + c_- \int_1^{\infty} e^{-i\rho|x| + it\rho} \rho^{-\lambda} (\rho|x|)^{-1/2} \rho^{d/2} d\rho \right) + |x|^{-\frac{d-2}{2}} \int_1^{\infty} e^{it\rho} \rho^{-\lambda} R_{\frac{d-2}{2}}(\rho|x|) \rho^{d/2} d\rho \nonumber \\
&= |x|^{-\frac{d-1}{2}} \left(c_+ \int_{1}^{\infty}  e^{i \rho ( t + |x|)}  \rho^{\frac{d-1}{2} - \lambda} d\rho + c_- \int_{1}^{\infty}  e^{i \rho( t - |x|)}  \rho^{\frac{d-1}{2} - \lambda} d\rho \right) \nonumber \\
&+ |x|^{-\frac{d-2}{2}} \int_{1}^{\infty} e^{it\rho} R_{\frac{d-2}{2}}(\rho |x|) \rho^{\frac{d}{2} -\lambda} d\rho   = |x|^{-\frac{d-1}{2}}(*) + |x|^{-\frac{d-2}{2}}(**).
\end{align*}
Assume that $\frac{d-1}{2} - \lambda \notin -\mathbb{N}$. Thus, for $(*)$ we can use Lemma \ref{lemma:invft-xalpha-nu} with $\alpha = \frac{d-1}{2} - \lambda$, and $\theta = t + |x|$ for the first term and $\theta = t - |x|$ for the second one:
\begin{align*}
(*) &= J\left(\frac{d-1}{2}-\lambda\right) \left\{c_+ [ e^{i\frac{\pi}{2}(\frac{d+1}{2} - \lambda)} (|x| + t)_+^{\lambda-\frac{d+1}{2}}\right. \\
&\hspace{48mm} + e^{-i\frac{\pi}{2}(\frac{d+1}{2} - \lambda)} \left(|x| + t\right)_-^{\lambda-\frac{d+1}{2}}] \nonumber \\
&\hspace{48mm} + c_-  [ e^{i\frac{\pi}{2}(\frac{d+1}{2} - \lambda)}(-|x| + t)_+^{\lambda-\frac{d+1}{2}} \\
&\left. \hspace{48mm}+ e^{-i\frac{\pi}{2}(\frac{d+1}{2} - \lambda)} \left( -|x| + t\right)_-^{\lambda-\frac{d+1}{2}} ] \right\}
\end{align*}
With respect to $(**)$, we consider the inequality \eqref{eq:ourbessel-asymp} to obtain that
\begin{align*}
&|x|^{-\frac{d-2}{2}}\left|  (**) \right| \lesssim |x|^{-\frac{d-1}{2}} \int_{1}^{\infty}\rho^{\frac{d-1}{2} - \lambda} \min\{ 1 , (\rho|x|)^{-1} \} d\rho \nonumber.
\end{align*}
If we bound the integrand by $\min\{ 1 , (\rho|x|)^{-1} \} \leq (\rho|x|)^{-1}$ we obtain
\begin{align*}
|x|^{-\frac{d-2}{2}}\left|  (**) \right| \lesssim \left( \int_{1}^{\infty} \rho^{\frac{d-3}{2} - \lambda} d\rho \right) |x|^{-\frac{d+1}{2}},
\end{align*}
where the integral in the right hand side converges if $\lambda \in (\frac{d-1}{2},d)$. On the other hand, if we bound the integrand by $\min\{ 1 , (\rho|x|)^{-1} \} \leq 1$, then we obtain
\begin{align*}
|x|^{-\frac{d-2}{2}}\left|  (**) \right| \lesssim \left( \int_1^{\infty} \rho^{\frac{d-1}{2} - \lambda} d\rho \right) |x|^{-\frac{d-1}{2}}.
\end{align*}
Now, the integral in the right hand side of the last inequality converges if $\lambda \in (\frac{d+1}{2},d)$. Denote $|x|^{-\frac{d-2}{2}} (**) = c_{\lambda}(x)$, so that 
\begin{align*}
|c_{\lambda}(x)| \lesssim \min\{|x|^{-\frac{d+1}{2}} , |x|^{-\frac{d-1}{2}}\}
\end{align*}
if $\lambda \in (\frac{d+1}{2},d)$, and 
\begin{align*}
|c_{\lambda}(x)| \lesssim |x|^{-\frac{d+1}{2}}
\end{align*}
if $\lambda \in (\frac{d-1}{2},d)$. This concludes the proof.
%Regarding $x = 0$, observe that
%\begin{align*}
%&\langle D \rangle^{\sigma} e^{\pm i t |D|} F_{\lambda}(0) = \int_{B(0,\nu)^c} e^{\pm i t |\xi|} |\xi|^{-\lambda} d\xi = w_{d-1} \int_{\mathbb{R}} \mathbb{I}_{(\nu,\infty)} \rho^{d-1-\lambda} e^{\pm i t \rho}  d\rho \nonumber \\
%&=w_{d-1} \mathcal{F}[\mathbb{I}_{(\nu,\infty)} \rho^{d-1-\lambda}](\pm t).
%\end{align*}
%By result \eqref{lemma:invft-xalpha-nu},
%\begin{align*}
%&\langle D \rangle^{\sigma} e^{\pm i t |D|} F_{\lambda}(0) =  w_{d-1}\{ \Gamma(d - \lambda) + J(n-\lambda-1,\nu) \}  \nonumber \\
%&\times \{ c_+ e^{i \frac{\pi}{2}(n-\lambda)} t_+^{-n + \lambda} +  c_- e^{-i \frac{\pi}{2}(n-\lambda)} t_-^{-n + \lambda} \}.
%\end{align*}
\end{proof}

\begin{proof}[Proof of Theorem \ref{thm:nonpropag-linear-rn}]

Let $A>1$. We already know from Theorem \eqref{thm:flambda-tnotzero} that $\langle D \rangle^{\sigma-1} e^{it|D|} F_{\lambda}$ has a finite norm in $L^p(B(0,A|t|)^c)$.\\

We focus on the part of the integration domain concerning $|x| < A|t|$. Namely, from Theorem \ref{thm:flambda-tnotzero} and its counterpart with $e^{-it|D|}$ we know that the key problematic parts of the integration domain to check whether 
$$\langle D \rangle^{\sigma-1} \sin(t|D|) F_{\lambda}$$
is in $L^p(\mathbb{R}^d)$ or not correspond to $\pm|x| \in (-t-\varepsilon,-t+\varepsilon)$ (regarding the $e^{it|D|}$ term) and $\pm|x| \in (t-\varepsilon,t+\varepsilon)$ (regarding the $e^{-it|D|}$ term), for $0 < \varepsilon \ll 1$. These points $x$ in $\mathbb{R}^d$ exist in $B(0,A|t|)$ because $A > 1$. By Theorem \ref{thm:flambda-tnotzero} and triangle inequality, the quantity
\begin{align*}
| \langle D \rangle^{\sigma-1} \sin(t|D|) F_{\lambda} \|_{L^p(\mathbb{R}^d)}
\end{align*}
is bounded below by the maximum between
\begin{align}\label{eq:boundbelow-1}
\left\| \right. & |x|^{-\frac{d-1}{2}} \left( c_+ [ e^{i\frac{\pi}{2}(\frac{d+1}{2} - \lambda)} (|x| + t)_+^{\lambda-\frac{d+1}{2}} + e^{-i\frac{\pi}{2}(\frac{d+1}{2} - \lambda)} \left(|x| + t\right)_-^{\lambda-\frac{d+1}{2}}] \nonumber  \right. \\
&\left. + c_-  [ e^{i\frac{\pi}{2}(\frac{d+1}{2} - \lambda)}(-|x| + t)_+^{\lambda-\frac{d+1}{2}} + e^{-i\frac{\pi}{2}(\frac{d+1}{2} - \lambda)} \left(-|x| + t\right)_-^{\lambda-\frac{d+1}{2}} ]\right. \nonumber \\
& \left. + c_+ [ e^{i\frac{\pi}{2}(\frac{d+1}{2} - \lambda)} (|x| - t)_+^{\lambda-\frac{d+1}{2}} + e^{-i\frac{\pi}{2}(\frac{d+1}{2} - \lambda)} \left( |x| - t \right)_-^{\lambda-\frac{d+1}{2}}] \nonumber \right. \\
&\left. \left. + c_-  [ e^{i\frac{\pi}{2}(\frac{d+1}{2} - \lambda)}(-|x| - t)_+^{\lambda-\frac{d+1}{2}} + e^{-i\frac{\pi}{2}(\frac{d+1}{2} - \lambda)} \left(-|x| - t\right)_-^{\lambda-\frac{d+1}{2}} ] \right) + c_{\lambda}  \right\|_{L^p(B(0,A|t|))} \nonumber \\
& \qquad - \left\| D_{\lambda} \right\|_{L^p(B(0,A|t|))}
\end{align}
and
\begin{align}\label{eq:boundbelow-2}
\left\| \right. & \left. |x|^{-\frac{d-1}{2}} \left( c_+ [ e^{i\frac{\pi}{2}(\frac{d+1}{2} - \lambda)} (|x| + t)_+^{\lambda-\frac{d+1}{2}} + e^{-i\frac{\pi}{2}(\frac{d+1}{2} - \lambda)} \left(|x| + t\right)_-^{\lambda-\frac{d+1}{2}}] \nonumber \right. \right. \\
&\left. + c_-  [ e^{i\frac{\pi}{2}(\frac{d+1}{2} - \lambda)}(-|x| + t)_+^{\lambda-\frac{d+1}{2}} + e^{-i\frac{\pi}{2}(\frac{d+1}{2} - \lambda)} \left(-|x| + t\right)_-^{\lambda-\frac{d+1}{2}} ]\right. \nonumber \\
& \left. + c_+ [ e^{i\frac{\pi}{2}(\frac{d+1}{2} - \lambda)} (|x| - t)_+^{\lambda-\frac{d+1}{2}} + e^{-i\frac{\pi}{2}(\frac{d+1}{2} - \lambda)} \left( |x| - t \right)_-^{\lambda-\frac{d+1}{2}}] \nonumber \right. \\
&\left. \left. + c_-  [ e^{i\frac{\pi}{2}(\frac{d+1}{2} - \lambda)}(-|x| - t)_+^{\lambda-\frac{d+1}{2}} + e^{-i\frac{\pi}{2}(\frac{d+1}{2} - \lambda)} \left(-|x| - t\right)_-^{\lambda-\frac{d+1}{2}} ]  \right) \right\|_{L^p(B(0,A|t|))} \nonumber \\
&\qquad - \left\| c_{\lambda} + D_{\lambda} \right\|_{L^p(B(0,A|t|))}.
\end{align}
Recall that $D_{\lambda} \in C^{\infty}(\mathbb{R}^d)$, so in particular $D_{\lambda} \in L^p_{\text{loc}}(\mathbb{R}^d)$. We focus on the remaining terms, and distinguish between the following cases:

\begin{itemize}
\item Let $\lambda \in (\frac{d+1}{2},d)$. By \eqref{eq:boundbelow-1}
\begin{align}\label{eq:gralint-p+2}
\| &\langle D \rangle^{\sigma-1} \sin(t|D|) F_{\lambda} \|_{L^p(\mathbb{R}^d)} \nonumber \\
&\qquad \gtrsim_{D_{\lambda}} \left(\int_{B(0,A|t|)} |x|^{-\frac{(d-1)p}{2}} \right. \nonumber \\
&\qquad \times\left| c_+ [ e^{i\frac{\pi}{2}(\frac{d+1}{2} - \lambda)} (|x| + t)_+^{\lambda-\frac{d+1}{2}} + e^{-i\frac{\pi}{2}(\frac{d+1}{2} - \lambda)} \left( |x| + t \right)_-^{\lambda-\frac{d+1}{2}}] \right. \nonumber \\
&\qquad \left. \left. + c_-  [ e^{i\frac{\pi}{2}(\frac{d+1}{2} - \lambda)}(-|x| + t)_+^{\lambda-\frac{d+1}{2}} + e^{-i\frac{\pi}{2}(\frac{d+1}{2} - \lambda)} \left( -|x| + t \right)_-^{\lambda-\frac{d+1}{2}} ]\right. \right. \nonumber \\
&\qquad \left. \left. + c_+ [ e^{i\frac{\pi}{2}(\frac{d+1}{2} - \lambda)} (|x| - t)_+^{\lambda-\frac{d+1}{2}} + e^{-i\frac{\pi}{2}(\frac{d+1}{2} - \lambda)} \left( |x| - t \right)_-^{\lambda-\frac{d+1}{2}}] \nonumber \right. \right. \\
&\qquad + c_-  [ e^{i\frac{\pi}{2}(\frac{d+1}{2} - \lambda)}(-|x| - t)_+^{\lambda-\frac{d+1}{2}} + e^{-i\frac{\pi}{2}(\frac{d+1}{2} - \lambda)} \left( -|x| - t \right)_-^{\lambda-\frac{d+1}{2}} ] \nonumber \\
&\hspace{40mm} \left. \left. + |x|^{\frac{d-1}{2}}c_{\lambda}(x) \right|^p  dx\right)^{1/p}.
\end{align}
If we had divergence around $|x| = \pm t$, this would follow from
\begin{align}\label{eq:divt}
\left( \lambda - \frac{d+1}{2} \right)p + 1 \leq 0 \Longleftrightarrow \lambda \leq \frac{d+1}{2} - \frac{1}{p}.
\end{align}
However, since $\lambda > d - \frac{d}{p}$, then \eqref{eq:divt} implies that $p<2$, so there is no divergence around $|x| = \pm t$ for $p>2$. On the other hand, the divergence around $x = 0$ would follow from
\begin{align*}
-p\left( \frac{d-1}{2} \right) + n \leq 0 \Longleftrightarrow p \geq \frac{2d}{d-1}.
\end{align*}
This covers the case $p>2$, for which divergence in \eqref{eq:gralint-p+2} comes from the factor $|x|^{-\frac{p(d-1)}{2}}$.\\

\item Let $\lambda \in (\frac{d-1}{2},d)$. By \eqref{eq:boundbelow-2}
\begin{align}\label{eq:gralint-p-2}
\| &\langle D \rangle^{\sigma-1} \sin(t|D|) F_{\lambda} \|_{L^p(\mathbb{R}^d)} \nonumber \\
&\qquad \gtrsim_{D_{\lambda}} \left( \int_{B(0,A|t|)} |x|^{-\frac{(d-1)p}{2}}\nonumber \right. \\
& \qquad \times \left| c_+ [ e^{i\frac{\pi}{2}(\frac{d+1}{2} - \lambda)} (|x| + t)_+^{\lambda-\frac{d+1}{2}} + e^{-i\frac{\pi}{2}(\frac{d+1}{2} - \lambda)} \left(|x| + t\right)_-^{\lambda-\frac{d+1}{2}}] \nonumber \right. \\
& \qquad \left. + c_-  [ e^{i\frac{\pi}{2}(\frac{d+1}{2} - \lambda)}(-|x| + t)_+^{\lambda-\frac{d+1}{2}} + e^{-i\frac{\pi}{2}(\frac{d+1}{2} - \lambda)} \left(-|x| + t\right)_-^{\lambda-\frac{d+1}{2}} ]\right. \nonumber \\
& \qquad \left. + c_+ [ e^{i\frac{\pi}{2}(\frac{d+1}{2} - \lambda)} (|x| - t)_+^{\lambda-\frac{d+1}{2}} + e^{-i\frac{\pi}{2}(\frac{d+1}{2} - \lambda)} \left( |x| - t \right)_-^{\lambda-\frac{d+1}{2}}] \nonumber \right. \\
&\qquad  \left. \left. + c_-  [ e^{i\frac{\pi}{2}(\frac{d+1}{2} - \lambda)}(-|x| - t)_+^{\lambda-\frac{d+1}{2}} + e^{-i\frac{\pi}{2}(\frac{d+1}{2} - \lambda)} \left(-|x| - t\right)_-^{\lambda-\frac{d+1}{2}} ] \right|^p  dx \right)^{1/p} \nonumber \\
&\hspace{30mm} - \left(\int_{B(0,A|t|)} |x|^{-\frac{p(d+1)}{2}} dx\right)^{1/p}.
\end{align}
Notice that the factors $|x|^{-\frac{(d+1)p}{2}}$ and $|x|^{-\frac{(d-1)p}{2}}$ will not diverge as long as
\begin{align*}
-p\left(\frac{d+1}{2} \right) + d > 0 \Longleftrightarrow p < \frac{2d}{d+1},
\end{align*}
while the divergence around $|x| = \pm t$ now follows from \eqref{eq:divt}. 
\end{itemize}
In summary, the divergence of the norm in $L^p(\mathbb{R}^d)$ follows from one of the following cases:
\begin{itemize}
\item $p < \frac{2d}{d+1}$ and $\frac{d-1}{2} < \lambda \leq \frac{d+1}{2} - \frac{1}{p}$, with divergence around $\pm |x| = t$, 
\item $p \geq \frac{2d}{d-1}$ and $\frac{d+1}{2} < \lambda$ with divergence around $x = 0$. Notice that  $\lambda > \frac{d+1}{2}$ is implied by $\lambda > d - \frac{d}{p}$, required for $F_{\lambda} \in W^{\sigma-1,p}(\mathbb{R}^d)$.
\end{itemize}
\end{proof}


\section{Deterministic counterexample for $S(t)$ on $\mathbb{T}^d$}
\label{sec:localization}
Our objective in this Section is to periodize $F_{\lambda}$ and $\langle D \rangle^{\sigma} \frac{\sin(t|D|)}{|D|} F_{\lambda}$, and to prove that the periodized counterparts serve also as counterexample for $x \in \mathbb{T}^d$. Namely, $f_{\lambda}$ is given by
\begin{align}\label{eq:flambdadef-1}
\sum_{k \in \mathbb{Z}^d} \langle D \rangle^{\sigma - 1} F_{\lambda}(x+k), \quad x \in \mathbb{T}^d,
\end{align}
which is an element in $W^{\sigma-1,p}(\mathbb{T}^d)$ for $\sigma \geq 1$ and the values of $p$ and $\lambda$ given in Lemma \ref{lemma:flambda-wsigmap-per}.
 %We restrict this part of the analysis to $p > 2$. 

We will use the periodization strategy used in \cite{W65}, based on the Poisson summation formula. We start giving the main result of this Section, which proof is presented below.

\begin{theorem}\label{thm:divergence-torus}
Let $\sigma \geq 1$, $d \geq 2$, $\lambda \in (\frac{d-1}{2},d)$ (satisfying $\frac{d-1}{2} - \lambda \notin -\mathbb{N}$) and $f_{\lambda}$ given in \eqref{eq:flambdadef}. Assume one of the following hypothesis:
\begin{itemize}
\item $p \geq \frac{2d}{d-1}$ and $\lambda > \frac{d+1}{2}$,
\item $p < \frac{2d}{d+1}$ and $\lambda \leq \frac{d+1}{2} - \frac{1}{p}$.
\end{itemize}
Then 
\begin{align*}
S(t)(0,f_{\lambda}) \notin W^{\sigma,p}(\mathbb{T}^d)
\end{align*}
for any $t \in (-1,1) \setminus \{0\}$ if $p \geq \frac{2d}{d-1}$, and for any $t \in (-1/2,1/2) \setminus \{0\}$ if $p < \frac{2d}{d+1}$.
\end{theorem}

\begin{remark}
Regarding Theorem \ref{thm:nonpropag-linear}, notice that the inequality $\lambda > \frac{d+1}{2}$ needed for the case $p \geq \frac{2d}{d-1}$ in Theorem \ref{thm:divergence-torus} is implied by the condition $\lambda > d - \frac{d}{p}$ required for $f_{\lambda} \in W^{\sigma-1,p}(\mathbb{T}^d)$.
\end{remark}
\begin{comment}
\begin{remark}
Once we prove Theorem \ref{thm:divergence-torus} for $t \in (0,1/3)$, by time reversibility it follows also for $t \in (-1/3,0)$.
\end{remark}
\end{comment}

We now provide the analogous result to Proposition \ref{prop:Flambda-inWsigmap} for $f_{\lambda}$. Recall  Definition \ref{def:g-lambda} and the Fourier coefficients given in \eqref{eq:flambdadef}.
\begin{lemma}\label{lemma:flambda-wsigmap-per}
Let $\sigma \geq 1$. Then
\begin{itemize}
\item If $1 \leq p < \infty$ and $\lambda \in (d - \frac{d}{p},d)$, then $f_{\lambda} \in W^{\sigma-1,p}(\mathbb{T}^d)$.
\item If $2 \leq p < \infty$ and $\lambda > d - \frac{d}{p}$, then $f_{\lambda} \in W^{\sigma-1,p}(\mathbb{T}^d)$.
\end{itemize}
\end{lemma}
\begin{proof}
From Proposition \ref{prop:Flambda-inWsigmap}, $F_{\lambda} \in W^{\sigma-1,1}(\mathbb{R}^d)$ as long as $\lambda \in (0,d)$. Therefore, the Fourier expansion of 
\begin{align*}
\sum_{k \in \mathbb{Z}^d} \langle D \rangle^{\sigma - 1} F_{\lambda}(x+k)
\end{align*}
is given by the coefficients (see Theorem 2.4 from Chapter 7 in \cite{steinweiss})
\begin{align*}
\widehat{f_{\lambda}}(k) =  h(k) |k|^{-\lambda}, \quad k \in \mathbb{Z}^3.
\end{align*}	
Moreover, using Lemma \ref{lemma:Flambda} we obtain that, for any $x \in \mathbb{T}^d \setminus \{0\}$ with at least one component not equal to $1$,
\begin{align*}
\sum_{k \in \mathbb{Z}^d} \langle D \rangle^{\sigma - 1} F_{\lambda}(x+k) = \gamma(d,\lambda) |x|^{-n+\lambda} + G_{\lambda,d}(x) + \sum_{k \in \mathbb{Z}^d \setminus \{0\}} \langle D \rangle^{\sigma - 1} F_{\lambda}(x+k),
\end{align*}
where the last sum is a uniformly convergent sum of functions in $C^{\infty}(\mathbb{T}^d)$. Therefore, the statement $f_{\lambda} \in W^{\sigma-1,p}(\mathbb{T}^d)$ follows from
\begin{align*}
\gamma(d,\lambda) |x|^{-n+\lambda} + G_{\lambda,d}(x) \in L^p(\mathbb{T}^d),
\end{align*}
and the latter is true if $\lambda > d - \frac{d}{p}$. This proves the first statement.

Regarding the second statement, let $p>2$. By Hausdorff-Young inequality we have that
\begin{align*}
\|f_{\lambda}\|_{W^{\sigma-1,p}(\mathbb{T}^d)} \lesssim \|(g_{\lambda}(k) e^{ikx})_{k \in \mathbb{Z}^d} \|_{\l^{p'}(\mathbb{Z}^d)} = \left( \sum_{k \in \mathbb{Z}^d} |g_{\lambda}(k)|^{p'} \right)^{\frac{1}{p'}},
\end{align*}
where the last sum converges if $\lambda > d - \frac{d}{p}$. This concludes the proof.
\end{proof}

\begin{remark}
As we did in Section \eqref{sec:linearcex}, we can simplify the following computations from the objects $\langle D \rangle^{\sigma} \frac{e^{it|D|}}{|D|} F_{\lambda}$ to  $\langle D \rangle^{\sigma-1} e^{it|D|} F_{\lambda}$, given that $\frac{|D|}{\langle D \rangle}$ is a bounded operator in $L^q(\mathbb{T}^d)$ for any $q \geq 1$, see Theorem $3.6$ from Chapter 7 in \cite{steinweiss} and the remark done above for the analogous situation in $L^q(\mathbb{R}^d)$.
\end{remark}

\begin{proof}[Proof of Theorem \ref{thm:divergence-torus}]
Let $t \neq 0$, $|t| < 1$. Then, there exists some $\gamma > 0$ sufficiently small such that $|t| < 1 / (2\gamma + 1)$. From Theorem \ref{thm:flambda-tnotzero} we have that $\langle D \rangle^{\sigma-1} e^{it|D|} F_{\lambda} \in L^1(\mathbb{R}^d)$ if $\lambda > \frac{d-1}{2}$. Thus, applying again Theorem 2.4 from Chapter 7 in \cite{steinweiss} as we similarly did in Lemma \ref{lemma:flambda-wsigmap-per}, we have that the $L^1(\mathbb{T}^d)$ norm of
\begin{align}\label{eq:flambdadef-t}
\sum_{k \in \mathbb{Z}^d} \langle D \rangle^{\sigma-1} e^{it|D|} F_{\lambda}(x+k)
\end{align}
converges and its Fourier coefficients are given by
\begin{align}\label{eq:periodization-obj2}
\mathcal{F}\left[ \sum_{l \in \mathbb{Z}^d} \langle D \rangle^{\sigma-1} e^{it|D|} F_{\lambda}(x+l) \right](k) = \langle k \rangle^{\sigma-1} e^{it|k|} g_{\lambda}(k), \quad k \in \mathbb{Z}^d.
\end{align} 
Recall that we want to prove the divergence of
\begin{align*}
\|S(t)(0,f_{\lambda})\|^p_{W^{\sigma,p}(\mathbb{T}^d)} = \frac{1}{2} \int_{\mathbb{T}^d} \left| \langle D \rangle^{\sigma} \frac{e^{it|D|} - e^{-it|D|}}{|D|} f_{\lambda}(x)\right|^p dx,
\end{align*}
which holds true if we prove the divergence of 
\begin{align}\label{eq:finaldiv-0}
\int_{\mathbb{T}^d} \left| \langle D \rangle^{\sigma-1} \sin(t|D|) f_{\lambda}(x)\right|^p dx.
\end{align}
Therefore, if we show that
\begin{align}\label{eq:finaldiv-2}
\sum_{k \in \mathbb{Z}^d} \langle D \rangle^{\sigma-1} \sin(t|D|) F_{\lambda}(x+k)
\end{align}
is not in $L^p(\mathbb{T}^d)$, the proof would be concluded. 

For a.e. $x \in \mathbb{T}^d \setminus \{0\}$ with at least one component not equal to $1$, and such that $|x+l| \neq t$ and $-|x+l| \neq t$ for any $l \in \mathbb{Z}^d$, from Theorem \ref{thm:flambda-tnotzero} we have that
\begin{align}\label{eq:finaldiv-3}
\langle D &\rangle^{\sigma-1} \sin(t|D|) F_{\lambda}(x+k)\nonumber \\
&= \frac{J\left(\frac{d-1}{2}-\lambda\right)}{2i} |x+k|^{-\frac{d-1}{2}} \nonumber \\
&\qquad \times  \{ c_{+} [ e^{i\frac{\pi}{2}(\frac{d+1}{2} - \lambda)} (|x+k| + t)_+^{\lambda-\frac{d+1}{2}} + e^{-i\frac{\pi}{2}(\frac{d+1}{2} - \lambda)} (|x+k| + t)_-^{\lambda-\frac{d+1}{2}}] \nonumber \\
&\qquad + c_{-}  [ e^{i\frac{\pi}{2}(\frac{d+1}{2} - \lambda)}(-|x+k| + t)_+^{\lambda-\frac{d+1}{2}} + e^{-i\frac{\pi}{2}(\frac{d+1}{2} - \lambda)} (-|x+k| + t)_-^{\lambda-\frac{d+1}{2}} ]\nonumber  \\
&\qquad - c_{+} [ e^{i\frac{\pi}{2}(\frac{d+1}{2} - \lambda)} (|x+k| - t)_+^{\lambda-\frac{d+1}{2}} + e^{-i\frac{\pi}{2}(\frac{d+1}{2} - \lambda)} (|x+k| - t)_-^{\lambda-\frac{d+1}{2}}] \nonumber \\
&\qquad - c_{-}  [ e^{i\frac{\pi}{2}(\frac{d+1}{2} - \lambda)}(-|x+k| - t)_+^{\lambda-\frac{d+1}{2}} + e^{-i\frac{\pi}{2}(\frac{d+1}{2} - \lambda)} (-|x+k| - t)_-^{\lambda-\frac{d+1}{2}} ] \}
\nonumber \\
&\qquad + c_{\lambda}(x+k) + D_{\lambda}(x+k)
\end{align}
where $k \in \mathbb{Z}^d$, $D_{\lambda} \in C^{\infty}(\mathbb{R}^d)$, $c_{\pm}$ the constants from \eqref{eq:asympBessel}, and 
\begin{align*}
&|c_{\lambda}(x+k)| \lesssim |x+k|^{-\frac{d+1}{2}} \text{ for } \lambda \in \left(\frac{d-1}{2},d\right), \nonumber \\
&|c_{\lambda}(x+k)| \lesssim |x+k|^{-\frac{d-1}{2}} \text{ for } \lambda \in \left(\frac{d+1}{2},d\right).
\end{align*}
From the proof of Theorem \ref{thm:nonpropag-linear-rn}, we expect that the singularity around $x = 0$ will yield the divergence for the case $p \geq \frac{2d}{d-1}$, while the singularity around $|x| = \pm |t|$ will give the divergence for $p < \frac{2d}{d+1}$. For the first case, by the triangle inequality,
\begin{align}\label{eq:trineq-tndiv}
&\left\| \sum_{k \in \mathbb{Z}^d}  \langle D \rangle^{\sigma-1} \sin(t|D|) F_{\lambda}(x+k)  \right\|_{L^p(\mathbb{T}^d)} \nonumber \\
&\qquad  \geq \left\| \sum_{k \in \mathbb{Z}^d} \langle D \rangle^{\sigma-1} \sin(t|D|) F_{\lambda}(x+k)  \right\|_{L^p(B(0,\gamma |t|))} \nonumber \\
&\qquad \geq \left\|  \langle D \rangle^{\sigma-1} \sin(t|D|) F_{\lambda}(x)  \right\|_{L^p(B(0,\gamma |t|))} \nonumber \\
&\qquad \qquad - \left\| \sum_{k \in \mathbb{Z}^d \setminus \{0\}} \langle D \rangle^{\sigma-1} \sin(t|D|) F_{\lambda}(x+k)  \right\|_{L^p(B(0,\gamma |t|))}.
\end{align}
Given $x \in B(0,\gamma |t|)$ and $k \in \mathbb{Z}^d \setminus \{0\}$, if we use $|t| < 1/(2\gamma+1)$,
\begin{align*}
&\left| |x+k| \pm t\right| \geq |x+k| - |t| \geq 1 - (\gamma + 1) |t|, \nonumber \\
& |x+k| \geq |k| - |x| \geq 1 - \gamma |t| > |t|,
\end{align*}
which are uniform and strictly positive lower bounds. Thus, regarding the second norm in the last expression of \eqref{eq:trineq-tndiv}, we use that for any $k \in \mathbb{Z}^d$ with $|k| > 1$ and $x \in \mathbb{T}^d$
\begin{align*}
|x + k| \geq |k| - |x| \geq |k| - 1,
\end{align*}
so that, by the second point in Theorem \ref{thm:flambda-tnotzero}, for any $N > \max\{d - \lambda,d/p\}$
\begin{align*}
&\left( \int_{B(0,\gamma |t|)} \left| \sum_{k \in \mathbb{Z}^d \setminus \{0\}} \langle D \rangle^{\sigma-1} \sin(t|D|) F_{\lambda}(x+k) \right|^p dx \right)^{1/p} \nonumber \\
&\lesssim \left(  \frac{1}{|t|^{Np}} + \sum_{k \in \mathbb{Z}^d, |k| > 1} \int_{B(0,\gamma |t|)} \frac{dx}{|x + k|^{Np}} \right)^{1/p} \nonumber \\
&\lesssim \left( \frac{1}{|t|^{Np}} + \sum_{k \in \mathbb{Z}^d, |k| > 1} \frac{1}{(|k| - 1)^{Np}}  \right)^{1/p} \lesssim 1.
\end{align*}
Therefore, the same analysis of the proof for Theorem \ref{thm:nonpropag-linear-rn} for the integral
\begin{align*}
\left( \int_{B(0,\gamma |t|)} \left| \langle D \rangle^{\sigma-1} \sin(t|D|) F_{\lambda}(x)  \right|^p dx \right)^{1/p}
\end{align*}
concludes the proof for $p \geq \frac{2d}{d-1}$. 

For $p < \frac{2d}{d+1}$, consider $|t| < 1/2$, $t \neq 0$. Analogously, there exists some $\gamma_1 > 0$ sufficiently small such that $|t| < 1/(2 + 2\gamma_1)$. Given $x \in B(0, (1  + \gamma_1) |t|) \setminus B(0, (1 - \gamma_1)|t|)$ and $k \in \mathbb{Z}^d \setminus \{0\}$,
\begin{align*}
&\left| |x+k| - t\right| \geq |x+k| - |t| \geq 1 - (2 + \gamma_1)|t|, \nonumber \\
& |x+k| \geq |k| - |x| \geq 1 - (1 + \gamma_1)|t| > |t|,
\end{align*}
which are again uniform and strictly positive lower bounds. Thus, by an analogous analysis than the one for $p \geq \frac{2d}{d-1}$, we have that
\begin{align*}
&\left\| \sum_{k \in \mathbb{Z}^d}  \langle D \rangle^{\sigma-1} \sin(t|D|) F_{\lambda}(x+k)  \right\|_{L^p(\mathbb{T}^d)} \nonumber \\
&\qquad \geq \left\|  \langle D \rangle^{\sigma-1} \sin(t|D|) F_{\lambda}(x)  \right\|_{L^p(B(0,(1 + \gamma_1)|t|) \setminus B(0,(1 - \gamma_1)|t|))} \nonumber \\
&\qquad \qquad - \left\| \sum_{k \in \mathbb{Z}^d \setminus \{0\}} \langle D \rangle^{\sigma-1} \sin(t|D|) F_{\lambda}(x+k)  \right\|_{L^p(B(0,(1 + \gamma_1)|t|) \setminus B(0,(1 - \gamma_1)|t|))},
\end{align*}
where the second norm is finite, and the first norm diverges by the same analysis of the proof for Theorem \ref{thm:nonpropag-linear-rn}.
\end{proof}

\section{Proof of Theorem \ref{thm:nonpropag-nonlinear}}\label{sec:proof-nonlinearcex}
Fix $d = 3$. All we need is to prove an appropriate control of \eqref{eq:quantity-to-control}. First of all, we identify the Sobolev spaces $H^{r-1}$ to which our initial velocity belongs.

\begin{lemma}\label{lemma:Hlambda-Hr}
	Let $\lambda \in (0,3)$ and $r < \sigma + \lambda - \frac{3}{2}$. Then $S_{\lambda} \subset H^{r-1}(\mathbb{T}^3)$. In particular, if $\lambda \in (\frac{3}{2},3)$, then $S_{\lambda} \subset H^{\sigma-1}(\mathbb{T}^3)$.
	\end{lemma}
	\begin{proof}
	We only need to check when $f_{\lambda} \in H^{r-1}(\mathbb{T}^3)$. This follows from
	\begin{align*}
	\|f_{\lambda}\|_{H^{r-1}(\mathbb{T}^3)}^2 = \sum_{k \in \mathbb{Z}^3} h(k) \langle k \rangle^{2(r-\sigma)} |k|^{-2\lambda},
	\end{align*}
	which converges if and only if $r < \sigma + \lambda - \frac{3}{2}$.
	\end{proof}
	
	
	To obtain a suitable control of the Duhamel integral for the first case in Theorem \ref{thm:nonpropag-linear} ($p \geq 3$), we use Strichartz estimates. In order to adapt them to the periodic setting, we use the finite speed of propagation of the linear flow from the linearization of the Cauchy system \eqref{eq:cauchywave} given for $x \in \mathbb{R}^d$. Denote by $\tilde{S}(t)$ the linear evolution from equation \eqref{eq:cauchywave} for $x \in \mathbb{R}^d$. This will allow us to give a periodic version of the well-known estimates on $\mathbb{R}^3$ from \cite{GV95}.
\begin{proposition}[finite speed of propagation]\label{prop:finite-speed-prop}
Let $F \in W^{\sigma-1,p}(\mathbb{R}^d)$ for some $\sigma \geq 1$ and $p \in [1,\infty)$ such that
\begin{align*}
\supp(F) \subset \{x \in \mathbb{R}^d : |x - x_0| \leq R\}
\end{align*}
for some $R > 0$ and $x_0 \in \mathbb{R}^d$. Then, for any $t \geq 0$,
\begin{align*}
\supp(\tilde{S}(t)(0,F)) \subset \{x \in \mathbb{R}^d : |x - x_0| \leq R + t\}.
\end{align*}
\end{proposition}
\begin{proof}
%It suffices to do a straightforward adaptation of the analogous result in \cite{Tzvetkov2019-notes} to the $W^{\sigma,p}(\mathbb{R}^d)$ setting.
We assume we know the result for $F \in C^1(\mathbb{R}^d)$, see Proposition 1.10 in \cite{Tzvetkov2019-notes}. Following this strategy, let $F \in W^{\sigma-1,p}(\mathbb{R}^d)$ and $\rho_{\varepsilon}(x)  = \varepsilon^{-n} \rho(x/\varepsilon)$, $\rho \in C^{\infty}_c(\mathbb{R}^d)$, $0 \leq \rho \leq 1$, and $\int_{\mathbb{R}^d} \rho(x) dx = 1$. Observe that
\begin{align*}
\rho_{\varepsilon} \ast \tilde{S}(t)(0,F) = \tilde{S}(t)(0,\rho_{\varepsilon} \ast F).
\end{align*}
Let $\tilde{\varphi} \in C^{\infty}_c(|x| > t + R)$. Then $\tilde{S}(t)(0,\rho_{\varepsilon} \ast F)(\tilde{\varphi}) = 0$ for $\varepsilon$ small enough. Also, $\rho_{\varepsilon} \ast \tilde{S}(t)(0,F)(\tilde{\varphi})$ converges to $\tilde{S}(t)(0,F)(\tilde{\varphi})$ as $\varepsilon \rightarrow 0^+$.
\end{proof}
Given $r \in \mathbb{R}$, $p,q \in [1,\infty]$, we denote 
\begin{align*}
L^q_T W^{r,p}(\mathbb{T}^3) = L^q([0,T] , W^{r,p}(\mathbb{T}^3)),
\end{align*}
being the external Lebesgue space given for the time variable, and the internal space for the space variable. 

\begin{proposition}\label{prop:st}
		Let $\mu \in \mathbb{R}$, $T \in (0,1)$ and $2 \leq p,q \leq \infty$ be such that
		\begin{align}\label{eq:adm-st}
			\frac{2}{q} \leq 1 - \frac{2}{p}, \quad (q,p) \neq (2,\infty), \quad \frac{3}{2} - \mu \leq \frac{3}{p} + \frac{1}{q}.
		\end{align} 
		We say that $(q,p)$ are $\mu$-admissible if they satisfy these relations. Then
		\begin{align*}
			\left\|\frac{\sin(t|D|)}{|D|} f\right\|_{L^q_TL^p(\mathbb{T}^3)} \lesssim \|g\|_{H^{\mu-1}(\mathbb{T}^3)}. 
		\end{align*}
	\end{proposition}

\begin{proof}
Let $\delta \in (0,1)$, and $(h_j)_{j \leq N_1}$ be a partition of unity in $\mathbb{T}^3$ such that 
\begin{align*}
\supp (h_j) \subset B(x_j,\delta), \quad j  = 1,\dots,N_{\delta},
\end{align*}
for certain $x_j \in \mathbb{T}^3$, $N_{\delta} \in \mathbb{N}$ and $h_j \in C^{\infty}(\mathbb{T}^3)$ for any $j \in \{1,\dots,N_{\delta}\}$. Moreover, for any $x \in \mathbb{T}^3$
\begin{align*}
\sum_{j = 1}^{N_{\delta}} h_j(x) = 1.
\end{align*} 
Then
\begin{align}
&\left\| \frac{\sin(t|D|)}{|D|} f \right\|_{L^q_TL^p(\mathbb{T}^3)} = \left\| \frac{\sin(t|D|)}{|D|} \left( \sum_{j=1}^{N_{\delta}} h_j \right)f \right\|_{L^q_TL^p(\mathbb{T}^3)} \nonumber \\
&\leq \sum_{j=1}^{N_{\delta}}  \left\| \frac{\sin(t|D|)}{|D|}\left(h_j f\right)\right\|_{L^q_TL^p(\mathbb{T}^3)}. \label{eq:loc-gral-1}
\end{align}
Recall that $t \in [0,T]$ for $T < 1$. Regarding each term in the last sum in \eqref{eq:loc-gral-1}, if we take $\delta \in (0,1)$  such that $\delta + t \leq 1$, then we can apply Strichartz estimates on $\mathbb{R}^3$ for any $\mu$-admissible pair $(q,p)$ thanks to the finite speed of propagation property from Proposition \ref{prop:finite-speed-prop}:
\begin{align*}
&\left\| \frac{\sin(t|D|)}{|D|} (h_jf) \right\|_{L^q_TL^p(\mathbb{T}^3)} = \left\| \frac{\sin(t|D|)}{|D|} (h_jf) \right\|_{L^q_TL^p(\mathbb{R}^3)} \lesssim \|h_j f\|_{H^{\mu-1}(\mathbb{R}^3)} \nonumber \\
&= \|h_j f\|_{H^{\mu-1}(\mathbb{T}^3)} \lesssim  \|h_j\|_{L^{\infty}(\mathbb{T}^3)} \|f\|_{H^{\mu-1}(\mathbb{T})^3)}  \lesssim_{h_j} \|f\|_{H^{\mu-1}(\mathbb{T})^3)} 
\end{align*}
This concludes the proof.
\end{proof}
	
	\begin{lemma}\label{lemma:infcontrol-duhamel-st}
	Let $T \in (0,1)$, $p \geq 2$, $\sigma > \frac{3}{p} - \frac{1}{2}$. Then
		\begin{align}\label{eq:control-duh}
			\left\| \int_0^t \frac{\sin((t-\tau)|D|)}{|D|} u^3(\tau) d\tau \right\|_{L^{\infty}_T W^{\sigma,p}(\mathbb{T}^3)} \lesssim T \sup_{|t| \leq T} \| u(t) \|_{L^{\infty}(\mathbb{T}^3)}^2 \|u(t)\|_{H^{\sigma + \frac{1}{2} - \frac{3}{p}}(\mathbb{T}^3)}.
		\end{align}
		Moreover
		\begin{align}\label{eq:control-duh-1}
		\left\| \int_0^t \frac{\sin((t-\tau)|D|)}{|D|} F(\tau) d\tau \right\|_{L^{\infty}_T W^{\sigma,p}(\mathbb{T}^3)} \lesssim  \|F\|_{L^1_T H^{\sigma + \frac{1}{2} - \frac{3}{p}}(\mathbb{T}^3)}.
		\end{align}
	\end{lemma}
	\begin{proof}
	Let $(q,p)$ be a $\mu$-admissible from \eqref{eq:adm-st}). Then, by the triangle and Minkowski inequalities,
	\begin{align*}
			&\left\| \int_0^t \frac{\sin((t-\tau)|D|)}{|D|} F(\tau) d\tau \right\|_{L^q_T W^{\sigma,p}(\mathbb{T}^3)}  \nonumber\\
			&\leq \left\| \left\| \left\|\frac{\sin((t-\tau)|D|)}{|D|} F(\tau)\right\|_{W^{\sigma,p}(\mathbb{T}^3)} \right\|_{L^1_{\tau}[0,T]} \right\|_{L^q_t[0,T]}  \nonumber\\
			&\leq \left\|  \left\|\frac{\sin((t-\tau)|D|)}{|D|} F(\tau)\right\|_{L^q_t([0,T],W^{\sigma,p}(\mathbb{T}^3))}  \right\|_{L^1_{\tau}[0,T]} \nonumber\\
			&\lesssim \left\| \left\| F(\tau) \right\|_{H^{\mu+\sigma-1}(\mathbb{R}^3)} \right\|_{L^1_{\tau}[0,T]} = \|F\|_{L^1_T H^{\sigma + \mu - 1}(\mathbb{T}^3)}
		\end{align*}
		Considering the sharp case in the last inequality of \eqref{eq:adm-st},
		\begin{align*}
		\mu = \frac{3}{2} - \frac{3}{p} - \frac{1}{q} \Longrightarrow \mu+\sigma-1 = \sigma + \frac{1}{2} - \frac{3}{p} - \frac{1}{q}.
		\end{align*}
		Take $q = \infty$ to obtain \eqref{eq:control-duh-1}. Furthermore, if we consider $\sigma + \frac{1}{2} - \frac{3}{p} > 0$ and $F = u^3$ we can apply Leibniz fractional rule \cite{grafakos_kp} to obtain
		\begin{align*}
		\|F\|_{L^1_T H^{\sigma + \frac{1}{2} - \frac{3}{p}}(\mathbb{T}^3)} \leq T \sup_{|t| \leq T} \|u(t)\|^2_{L^{\infty}(\mathbb{T}^3)} \|u(t)\|_{H^{\sigma + \frac{1}{2} - \frac{3}{p}}(\mathbb{T}^3)}.
		\end{align*}
	\end{proof}
	
	
	\begin{proof}[Proof of Theorem \ref{thm:nonpropag-nonlinear}]
	Recall that $p \geq 3$, $\lambda \in (3-\frac{3}{p},3)$ and $\sigma > 1$, and notice that under these conditions we cannot have $1 - \lambda \in -\mathbb{N}$.
	
	The density of $S_{\lambda}$ in $W^{\sigma - 1,p}(\mathbb{T}^3)$ follows from the density of $C^{\infty}(\mathbb{T}^3)$. Furthermore, since by Theorem \ref{thm:nonpropag-linear} we have
	\begin{align*}
	\frac{\sin(t|D|)}{|D|} f_{\lambda} \notin W^{\sigma,p}(\mathbb{T}^3)
	\end{align*}
	for any $p \geq 3$, $\sigma \geq 1$ and $t \in (-1,1) \setminus \{0\}$, then we only need to prove
	\begin{align}\label{eq:mainproof-2}
	\left\| \int_0^t \frac{\sin((t-\tau)|D|)}{|D|} \right. & \left. u^3(\tau) d\tau \right\|_{L^{\infty}_T W^{\sigma,p}(\mathbb{T}^3)} < \infty
	\end{align}
	for $T \in (0,1)$. To be more precise, once we prove \eqref{eq:mainproof-2}, then we would have shown that, for any $t \in (-1,1) \setminus \{0\}$,
	\begin{align*}
	u(\cdot,t) \notin W^{\sigma,p}(\mathbb{T}^3).
	\end{align*}
	Our choices of $\lambda$, $p$ and $\sigma$ let us bound the Duhamel integral in two steps:
	
	\begin{itemize}
	\item From Lemma \ref{lemma:Hlambda-Hr}, $S_{\lambda} \subset H^{(\sigma  + \lambda - \frac{3}{2}-\varepsilon)-1}(\mathbb{T}^3)$ for any $\varepsilon > 0$. Observe that, with our choice of $\lambda$ and $p$,
	\begin{align*}
	\sigma + \lambda - \frac{3}{2} > \sigma + 3 \left( \frac{1}{2} - \frac{1}{p}\right) \geq \sigma + \frac{1}{2}.
	\end{align*} 
	Therefore $\sigma \geq \frac{1}{2}$ suffices to have existence of a unique global solution $u$ of \eqref{eq:cauchywave} in the space $C(\mathbb{R},H^r(\mathbb{T}^3))$ for any $r \in [1,\sigma + \lambda - \frac{3}{2})$ (see \cite{Tzvetkov2019-notes}). Namely, $u \in C(\mathbb{R},H^{\sigma + \frac{1}{2}}(\mathbb{T}^3))$.
	
	\item Let $T \in (0,1)$. By Lemma \ref{lemma:infcontrol-duhamel-st} we have that
	\begin{align}\label{eq:mainproof-1}
	\left\| \int_0^t \frac{\sin((t-\tau)|D|)}{|D|} \right. & \left. u^3(\tau) d\tau \right\|_{L^{\infty}_T W^{\sigma,p}(\mathbb{T}^3)} \nonumber \\
	&\lesssim T  \| u \|_{L^{\infty}_TL^{\infty}(\mathbb{T}^3)}^2 \|u\|_{L^{\infty}_T H^{\sigma  + \frac{1}{2} - \frac{3}{p}}(\mathbb{T}^3)},
	\end{align}
	where we used $\sigma > 1 \geq \frac{3}{p} - \frac{1}{2}$ for any $p \geq 3$ (this still holds for any $\sigma > \frac{1}{2}$). Notice that, from the previous point,
	\begin{align*}
	\|u\|_{L^{\infty}_T H^{\sigma  + \frac{1}{2} - \frac{3}{p}}(\mathbb{T}^3)} \leq \|u\|_{L^{\infty}_T H^{\sigma  + \frac{1}{2}}(\mathbb{T}^3)} < \infty.
	\end{align*}			
	Furthermore, given $\sigma > 1$ and the Sobolev embedding of $H^{\frac{3}{2} + \varepsilon}(\mathbb{T}^3)$ in $L^{\infty}(\mathbb{T}^3)$ for any $\varepsilon > 0$, we have that
	\begin{align*}
	\| u \|_{L^{\infty}_TL^{\infty}(\mathbb{T}^3)} \lesssim \| u \|_{L^{\infty}_T H^{\sigma + \frac{1}{2}}(\mathbb{T}^3)} < \infty.
	\end{align*}
	\end{itemize}
	\end{proof}
	
	\begin{remark}
	Consider the following estimates from \cite{SSS91}, \cite{Peral1980} (see also chapter $9$ of \cite{steinha}), suitably adapted to the periodic setting following the strategy from Proposition \ref{prop:st}:
	\begin{align}\label{eq:estimates_linear_torus}
		&\left\| \frac{\sin(t|D|)}{|D|} f \right\|_{W^{\alpha,p}(\mathbb{T}^d)} \lesssim \| f \|_{L^p(\mathbb{T}^d)}, \quad t \in (0,1),
	\end{align}
	with
	\begin{align*}
	&\left| \frac{1}{p} - \frac{1}{2} \right| \leq \frac{1}{d-1}, \quad \alpha \leq 1 - (d-1)\left| \frac{1}{p} - \frac{1}{2}\right|, \quad d \geq 2.
	\end{align*}
	These estimates also allow us to control \eqref{eq:quantity-to-control} when $p \geq 3$. This comes from the following inequality, given $T \in (0,1)$:
	\begin{align*}
			\left\| \int_0^t \frac{\sin((t-\tau)|D|)}{|D|} u^3(\tau) d\tau \right\|_{L^{\infty}_T W^{\sigma,p}(\mathbb{T}^3)} \lesssim T \| u \|_{L^{\infty}_T L^{\infty}(\mathbb{T}^3)}^2 \|u\|_{L^{\infty}_T H^{\sigma + \frac{3}{2} - \frac{5}{p}}(\mathbb{T}^3)}.
		\end{align*}	
	
	Indeed, the integrand in $\tau$ of the Duhamel integral \eqref{eq:quantity-to-control}, given by $v(x,t,\tau) = \frac{\sin((t-\tau)|D|)}{|D|} u^3(x,\tau)$, solves the problem
	\begin{align*}
    \left\{
			\begin{array}{l}
				\partial_{tt} v - \Delta v = 0, \quad x \in \mathbb{T}^3, \quad
				\\
				v(x,\tau,\tau) = 0,  \ \ \partial_t v(x,\tau,\tau) = u^3(x,\tau).
			\end{array}
			\right.
	\end{align*}
	Therefore, from \eqref{eq:estimates_linear_torus},
	\begin{align*}
	\|v(\cdot,t,\tau)\|_{W^{\sigma,p}(\mathbb{T}^3)} \lesssim \|u^3(\cdot,\tau)\|_{W^{\sigma + \left(1 - \frac{2}{p}\right) - 1,p}}.
	\end{align*}
	Thus, by the fractional Leibniz rule,
	\begin{align}\label{eq:notmainproof-1}
	\left\| \int_0^t \frac{\sin((t-\tau)|D|)}{|D|} \right. & \left. u^3(\tau) d\tau \right\|_{L^{\infty}_T W^{\sigma,p}(\mathbb{T}^3)} \lesssim \sup_{|t| \leq T} \int_0^T \| u^3(\cdot,\tau)\|_{W^{\sigma - \frac{2}{p},p}(\mathbb{T}^3)}  d\tau \nonumber \\
	&\leq T \| u \|^2_{L^{\infty}_T L^{\infty}(\mathbb{T}^3)}\| u\|_{L^{\infty}_T W^{\sigma - \frac{2}{p},p}(\mathbb{T}^3)}.
	\end{align}
	Observe that 
\begin{itemize}
\item we have the Sobolev embedding $H^{\sigma + \frac{3}{2} - \frac{5}{p}}(\mathbb{T}^3) \hookrightarrow W^{\sigma - \frac{2}{p},p}(\mathbb{T}^3)$,
\item and $S_{\lambda} \subset H^{(\sigma + \frac{3}{2} - \frac{5}{p}) - 1}(\mathbb{T}^3)$ as long as $\lambda > 3 - \frac{5}{p}$ by Lemma \ref{lemma:Hlambda-Hr}.
\end{itemize}	
Therefore, given $p$, $\lambda$ and $\sigma$ as in Theorem \ref{thm:nonpropag-nonlinear}, in \eqref{eq:notmainproof-1} we can bound as
\begin{align*}
\left\| \int_0^t \frac{\sin((t-\tau)|D|)}{|D|} \right. & \left. u^3(\tau) d\tau \right\|_{L^{\infty}_T W^{\sigma,p}(\mathbb{T}^3)} \nonumber \\
	&\lesssim T \| u \|^2_{L^{\infty}_T H^{\sigma + \frac{1}{2}}(\mathbb{T}^3)}\| u\|_{L^{\infty}_T H^{\sigma + \lambda - \frac{3}{2} - \varepsilon}(\mathbb{T}^3)}
\end{align*}
for some $\varepsilon > 0$. Thus, we can proceed as in the proof of Theorem \ref{thm:nonpropag-nonlinear} to get the same result. The reason of having used Strichartz estimates instead of this procedure is that, by using estimates \eqref{eq:estimates_linear_torus}, we eventually require an $L^2$-based Sobolev regularity $\sigma + \frac{3}{2} - \frac{5}{p}$ to bound the left-hand side of \eqref{eq:notmainproof-1}, which for any $p \geq 2$ is larger than the analogous Sobolev regularity $\sigma + \frac{1}{2} - \frac{3}{p}$ obtained with Strichartz estimates (see Lemma \ref{lemma:infcontrol-duhamel-st}).
	\end{remark}
	
	\appendix
	
		\section{Appendix: Bessel functions}\label{app:bessel}
We define Bessel functions in the following way. See \cite{watson} or Chapter $4$ of \cite{steinweiss} for further detail. %Note that the lower bound on $k$ ensures the convergence of the integral in \eqref{eq:bessel2}.
	\begin{definition}\label{lemma:bessel_kreal}
		Let $k > -1/2$ and $t > 0$. Then
		\begin{align}\label{eq:bessel2}
			J_k(t) = \frac{(t/2)^k}{\Gamma(2k+1)\Gamma(1/2)} \int_{-1}^1 e^{its} (1-s^2)^{(2k-1)/2}ds.
		\end{align}
	\end{definition}
	
Moreover, for $t$ large we have the following asymptotic behaviour:
		\begin{align}\label{eq:asympBessel}
			J_k(t) = t^{-1/2} e^{it} \sum_{j=0}^N a_j t^{-j} + t^{-1/2} e^{-it} \sum_{j=0}^N b_j t^{-j} + R_k(t)
		\end{align}
for suitable $(a_j)_{j \geq 0}, (b_j)_{j \geq 0} \subset \mathbb{C}$,	for any $N \in \mathbb{N} \cup \{0\}$ and $R_k(t) = \mathcal{O}(t^{-N-\frac{3}{2}})$, where $\mathcal{O}$ is given for $t \rightarrow \infty$. 
\\


We present the proof of Lemma \ref{lemma:ftransformradial_bessel}, which justifies the application of Bessel functions in this work. 
	\begin{proof}[Proof of Lemma \ref{lemma:ftransformradial_bessel}]
	The case $f \in L^1(\mathbb{R}^d)$ is proved in page 155 of \cite{steinweiss}. To extend this result to any $f$ in $L^p(\mathbb{R}^d)$, $1 < p < \infty$, it suffices to use a density argument with radial mollifiers. Indeed, let $\phi \in C^{\infty}_c(\mathbb{R}^d)$ be radial and denote $\phi_{\varepsilon} = \varepsilon^{-n}\phi(\cdot/\varepsilon)$ for any $\varepsilon > 0$, so that $f_{\varepsilon} = \phi_{\varepsilon} \ast f \rightarrow f$ as $\varepsilon \rightarrow 0^+$ in $L^p(\mathbb{R}^d)$. In particular, $f_{\varepsilon}(x) \rightarrow f(x)$ for a.e. $x \in \mathbb{R}^d$ up to a subsequence. Observe that $f_{\varepsilon}$ is radial because $f$ and $\phi_{\varepsilon}$ are radial, so we can write $f_{\varepsilon}(x) = f_{\varepsilon,0}(|x|)$ for some $f_{\varepsilon,0}$, and apply the first part of the proof to each $\mathcal{F}[f_{\varepsilon}]$. Thus, $f_{\varepsilon,0}(|x|) \rightarrow f_0(|x|)$ for a.e. $x \in \mathbb{R}^d$ and $\varepsilon \rightarrow 0^+$, up to a subsequence.
	
Thanks to the dominated convergence theorem, for any $\varphi \in \mathcal{S}(\mathbb{R}^d)$ we have that
	\begin{align*}
	\int_{\mathbb{R}^d} f_{\varepsilon}(x) \varphi(x) dx \rightarrow \int_{\mathbb{R}^d} f(x) \varphi(x) dx \text{ as } \varepsilon \rightarrow 0^+.
	\end{align*}
	Moreover, 
	\begin{align*}
	&\int_{\mathbb{R}^d} f_{\varepsilon}(x) \varphi(x) dx = \int_{\mathbb{R}^d} \mathcal{F}[f_{\varepsilon}](\xi) \mathcal{F}[\varphi](\xi) d\xi \nonumber \\
	& \int_{\mathbb{R}^d} \left( c(d) |\xi|^{-\frac{d-2}{2}} \int_0^{\infty} f_{\varepsilon,0}(\rho) \rho^{\frac{d}{2}} J_{\frac{d-2}{2}}(\rho|\xi|) d\rho \right) \mathcal{F}[\varphi](\xi) d\xi.
	\end{align*}
	where we used Plancherel theorem and the known result in $L^1(\mathbb{R}^d)$. Since the integral
	\begin{align*}
	c(d) |\xi|^{-\frac{d-2}{2}} \int_0^{\infty} f_0(\rho) \rho^{\frac{d}{2}} J_{\frac{d-2}{2}}(\rho|\xi|) d\rho
	\end{align*}
	exists for a.e. $\xi \in \mathbb{R}^d$, by the a.e. pointwise convergence given above $f_{\varepsilon,0}(\rho) \rightarrow f_0(\rho)$ and by dominated convergence theorem we have that
	\begin{align*}
	\mathcal{F}[f](\xi) = c(d) |\xi|^{-\frac{d-2}{2}} \int_0^{\infty} f_0(\rho) \rho^{\frac{d}{2}} J_{\frac{d-2}{2}}(\rho|\xi|) d\rho
	\end{align*}
	for a.e. $\xi \in \mathbb{R}^d$.
	\end{proof}

	
	\section{Appendix: Homogeneous distributions}\label{app:temp-dist-xalpha+}
Let $\alpha \in \mathbb{C}$ and $\varphi \in \mathcal{S}(\mathbb{R})$. Observe that
\begin{align*}
		\int_0^{\infty} \mathbb{I}_{(1,\infty)}(x) x^{\alpha} \varphi(x) dx
	\end{align*}
	is well-defined for any complex value of $\alpha$. In this section we will compute the inverse distributional Fourier transform of $\mathbb{I}_{(1,\infty)}(x) x^{\alpha}$, i.e. we will prove Lemma \ref{lemma:invft-xalpha-nu}. We need a preliminary result.

\begin{lemma}\label{lemma:limit-complex}
	Let $\theta, \beta \in \mathbb{R}$. Then
	\begin{align*}
	\lim_{\gamma \rightarrow 0^+} (\theta + i \gamma)^{\beta} = \left\{ \begin{array}{lcc} \theta^{\beta} & \text{for} & \theta > 0 \\ \\ (-1)^{\beta}\theta^{\beta} e^{i\beta\pi} & \text{for} & \theta < 0. \end{array} \right.
	\end{align*}
	\end{lemma}
	
\begin{proof}
	Observe that
	\begin{align*}
		\lim_{\tau \rightarrow 0^+} &(\theta + i \tau)^{\beta} = \exp[\beta(\log|\theta + i\tau| + i \text{arg}(\theta+i\tau))] \nonumber\\
		&= \left\{ \begin{array}{lcc} \theta^{\beta} & \text{for} & \theta > 0 \\ \\ (-1)^{\beta}\theta^{\beta} e^{i\beta \pi} & \text{for} & \theta < 0, \end{array} \right. 
	\end{align*}
	given that $\text{arg}(\theta + i\tau) \overset{\tau \rightarrow 0^+}{\rightarrow} 0$ for $\theta > 0$ and $\text{arg}(\theta + i\tau) \overset{\tau \rightarrow 0^+}{\rightarrow} \pi$ for $\theta < 0$.
	\end{proof}
	
	\begin{proof}[Proof of Lemma \ref{lemma:invft-xalpha-nu}]
	Consider
	\begin{align*}
		G_{\tau}(x) = x^{\alpha} e^{-\tau x} \mathbb{I}_{(1,\infty)}(x),
	\end{align*}
	for $\tau > 0$, so that we will compute $\mathcal{F}^{-1}[G_{\tau}]$ and take the limit $\tau \rightarrow 0^+$ in $\mathcal{S}'(\mathbb{R})$. Given $\theta \in \mathbb{R}$, write
	\begin{align*}
		\mathcal{F}^{-1}[G_{\tau}](\theta) = \int_{1}^{\infty} x^{\alpha} e^{-\tau x} e^{ix\theta} dx = \int_{1}^{\infty} x^{\alpha} e^{i z x} dx, 
	\end{align*}
	for $z = \theta + i \tau$. This implies that $z$ is in the upper half-plane of $\mathbb{C}$, so that $\text{arg}(z) \in (0,\pi)$. We change variables through the rotation $\xi = - i z x$, so $\text{arg}(\xi) = \text{arg}(z)  - \frac{\pi}{2} \in (-\pi/2,\pi/2)$ and
	\begin{align*}						
		\int_{1}^{\infty} x^{\alpha} e^{i z x} dx = \left( \frac{e^{i\pi/2}}{z} \right)^{\alpha+1} \int_{e^{i(\text{arg}(z)-\pi/2)}(1,\infty)} \xi^{\alpha} e^{-\xi} d\xi.
	\end{align*}
	The next step is to prove that we can remove the rotation $e^{i(\text{arg}(z)-\pi/2)}$ from the integration domain. Consider within the complex plane a contour $L$ surrounding the sector between angles $0$ and $\text{arg}(z)-\pi/2$. This contour departs from radius $1$ and reaches some radius $M>0$ which we will make arbitrarily large eventually. 

Assume without loss of generality that $\text{arg}(z) \in (\frac{\pi}{2},\pi)$, the case $\text{arg}(z) \in (0,\frac{\pi}{2})$ is proved in a similar way. We write $L = L_1 + L_2 + L_3 + L_4$, where
		\begin{align*}
		&L_1 = \{ x : x \in (1,M) \}, \quad L_2 = \{ e^{i(\text{arg}(z)-\pi/2)}(M + 1 - x) : x \in (1,M) \}, \nonumber \\
			&L_3 = \{ M e^{i\gamma} : \gamma \in (0,\text{arg}(z)-\pi/2)\}, \nonumber \\
			& L_4 = \{e^{i(\text{arg}(z) - \pi/2-\gamma)} : \gamma \in (0,\text{arg}(z)-\pi/2)\}.
	\end{align*}
Given that $\xi^{\alpha} e^{-\xi}$ is holomorphic in $\{\xi \in \mathbb C : |\xi| > 0\}$, by Cauchy's theorem we have that
	\begin{align*}
		\int_L \xi^{\alpha} e^{-\xi} d\xi = 0.
	\end{align*}
	The term from that integral corresponding to a constant radius $M$ satisfies
	\begin{align*}
		&\left|\int_{L_3}\xi^{\alpha} e^{-\xi} d\xi \right| = \left| \int_0^{\text{arg}(z)-\pi/2} M^{\alpha+1} e^{i (\alpha + 1) \gamma} e^{-M(\cos \gamma + i \sin \gamma)} i d\gamma \right|  \nonumber\\
		&\leq M^{\text{Re}(\alpha) + 1} \left( \text{arg}(z)-\frac{\pi}{2} \right)\sup_{\gamma \in (0,\text{arg}(z)-\pi/2)} e^{- M \cos(\gamma)},
	\end{align*}
	which tends to $0$ as $M \rightarrow \infty$.
	 
	On the other hand, the sector of constant radius $1$ satisfies
	\begin{align*}
		&\left|\int_{L_4}\xi^{\alpha} e^{-\xi} d\xi \right| \nonumber \\
		&= \left| \int_0^{\text{arg}(z) - \pi/2}  e^{-i (\alpha + 1) \gamma} e^{-(\cos(\text{arg}(z) - \pi/2-\gamma) + i \sin(\text{arg}(z) - \pi/2-\gamma))} i d\gamma \right|  \nonumber\\
		&\leq  \left( \text{arg}(z) - \frac{\pi}{2} \right) \sup_{\gamma \in (0,\text{arg}(z)-\frac{\pi}{2})} e^{- \cos(\text{arg}(z) - \pi/2-\gamma)} \leq \frac{\pi}{2}.
	\end{align*}
	Therefore,
	\begin{align*}
	\mathcal{F}^{-1}[G_{\tau}](\theta) = \left(\int_0^{\infty} \xi^{\alpha} e^{-\xi} d\xi - \int_0^1 \xi^{\alpha} e^{-\xi} d\xi - \int_{L_4}\xi^{\alpha} e^{-\xi} d\xi \right)\left( \frac{e^{i\pi/2}}{\theta + i \tau} \right)^{\alpha+1}.
	\end{align*}
	Observe that
\begin{align*}
\max\left\{ \left|\int_0^{\infty} \right. \right. & \left. \left. e^{-\xi} \xi^{\alpha} d\xi \right| , \left|\int_0^{1} e^{-\xi} \xi^{\alpha} d\xi \right| \right\} \\
&\leq \int_0^{\infty} e^{-\xi} \xi^{\text{Re}(\alpha)} d\xi \leq \Gamma(\text{Re}(\alpha) + 1),
\end{align*}
where the latter is well-defined as long as we consider $\text{Re}(\alpha) \notin -\mathbb{N}$ and an analytic continuation of $\Gamma$ (see e.g. \cite{steinshak-complex}).
All in all,
\begin{align}\label{eq:ftransform-x+exp-translation}
\mathcal{F}^{-1}[G_{\tau}](\theta) = J(\alpha)  \left( \frac{e^{i\pi/2}}{\theta + i \tau} \right)^{\alpha+1},
\end{align}
where $J(\alpha)$ is a constant only depending on $\alpha$ such that
\begin{align*}
|J(\alpha)| \leq  \Gamma(\text{Re}(\alpha)+1) + \frac{\pi}{2}.
\end{align*}
If we take the limit $\tau \rightarrow 0^+$ in $\mathcal{S}'(\mathbb{R})$, from Lemma \ref{lemma:limit-complex} we have that
\begin{align}\label{eq:ftransform_x+-translation}
&\mathcal{F}^{-1}[x^{\alpha}\mathbb{I}_{(1,\infty)}(x)](\theta) \nonumber \\
&= J(\alpha)  [e^{i\pi(\alpha+1)/2} (\theta \mathbb{I}_{(0,\infty)}(\theta))^{-\alpha-1} + e^{-i\pi(\alpha+1)/2} (|\theta| \mathbb{I}_{(-\infty,0)}(\theta))^{-\alpha-1}].
\end{align}
\end{proof}

\section*{Acknowledgments}
	The author is supported by the Basque Government through the program BERC 2026-2029 (BCAM), by the project PID2024-156169NB-I00 (NCAFA), by the Severo Ochoa accreditation CEX2021-001142-S (BCAM), and by the predoctoral program of the Education Department of the Basque Government.
	
	The author thanks Nikolay Tzvetkov for proposing this interesting problem during his stay at L'École Normale Supérieure de Lyon, as well as his PhD supervisor Renato Lucà for useful advices.

	%\nocite{*}
	%\printbibliography
	
	
		\begin{thebibliography}{10}
	
		\bibitem{AT08}
		A.~Ayache and N.~Tzvetkov.
		\newblock The $L^p$ properties of Gaussian random series.
		\newblock {\em Trans. Am. Math.}, 360(8):4425–-4439, 2008.

		
		
		%\bibitem{B82}
		%M.~Beals.
		%\newblock $L^p$ Boundedness of Fourier Integral Operators.
		%\newblock {\em Mem. Amer. Math. Soc.}, 264, 1982.	
		
		\bibitem{Bogachev1998}
		V.I.~Bogachev.
		\newblock {\em Gaussian measures}.
		\newblock Mathematical surveys and monographs, 62 (American Mathematical Society, Providence, RI, 1998).
    	
		\bibitem{BurqTzvet}
		N.~Burq and N.~Tzvetkov.
		\newblock Random data {C}auchy theory for supercritical wave equations. {I}.
		{L}ocal theory.
		\newblock {\em Invent. Math.}, 173(3):449--475, 2008.
		
		\bibitem{CG23}
		N.~Camps and L.~Gassot.
		\newblock Pathological set of initial data for scaling-supercritical nonlinear Schrödinger equations.
		{L}ocal theory.
		\newblock {\em Int. Math. Res. Not.}, 2023(15):13214--13254, 2023.
		
		\bibitem{PZ14}
		G.~Da Prato and J.~Zabczyk.
		\newblock {\em Stochastic Equations in Infinite Dimensions}, 				2nd Edition (Cambridge University Press, 2014).
		
		\bibitem{fernique}
		X.~Fernique.
		\newblock Regularité des trajectoires des fonctions aléatoires Gausiennes.
		\newblock {\em Ecole d’Eté de Probabilités de St. Flour IV-1974. Lecture Notes in Mathematics}, 480, 1--96, 1975.
		
		\bibitem{GS64}
		I.M.~Gelfand and G.E.~Shilov.
		\newblock {\em Generalized functions, Vol. I: Properties and Operations}.
		\newblock Academic Press (American Mathematical Society, Providence, RI, 1964).
		

		\bibitem{GV95}
		J.~Ginibre and G.~Velo.
		\newblock Generalized Strichartz Inequalities for the Wave Equation.
		\newblock {\em Partial Differential Operators and Mathematical Physics. Oper. Theory Adv. Appl.}, 78, 1995.
		
		\bibitem{grafakos_kp}
		L.~Grafakos and S.~Oh.
		\newblock The Kato-Ponce inequality.
		\newblock {\em Commun. Partial Differ. Equ.}, 39(6):1128–-1157, 2014.	
		
		%\bibitem{GG15}
		%P.~Gérard and S.~Grellier.
		%\newblock The cubic Szëgo equation and Hankel operators.
		%\newblock {\em arXiv preprint arXiv:1508.06814}, 2015.
		
		%\bibitem{GG17}
		%P.~Gérard and S.~Grellier.
		%\newblock Generic colourful tori and inverse spectral transform for Hankel operators.
		%\newblock {\em Tunisian Journal of Mathematics}, 1(3):347--372, 2018.
		
		\bibitem{GOTW22}
		T.~Gunaratnam, T.~Oh, N.~Tzvetkov and H.~Weber.
		\newblock Quasi-invariant Gaussian measures for the nonlinear wave equation in three dimensions.
		\newblock {\em Probab. Math. Phys.}, 3(2):343–-379, 2022.
		
		\bibitem{Miyachi1980}
		A.~Miyachi.
		\newblock On some estimates for the wave equation in $L^p$ and $H^p$.
		\newblock {\em J. Fac. Sci. Univ. Tokyo Sect. IA Math.}, 27, 1980.		
		
		
		%\bibitem{Ha14}
		%Z.~Hani.
		%\newblock Long-time instability and unbounded Sobolev orbits for some periodic nonlinear Schrödinger equations.
		%\newblock {\em Archive for Rational Mechanics and Analysis}, 211(3):929---964, 2014.
	
	
		%\bibitem{HNVW16}
		%T.~Hytönen, J.~ van Neerven, M.~Veraar and L.~Weis.
		%\newblock {\em Analysis in Banach Spaces Volume I: Martingales and Littlewood-Paley Theory}.
		%\newblock Ergebnisse der Mathematik und ihrer Grenzgebiete. 3. Folge / A Series of Modern Surveys in Mathematics. Springer International, 2016.
		
		\bibitem{MS1}
		C.~Muscalu, W.~ Schlag.
		\newblock {\em Classical and Multilinear Harmonic Analysis}.
		\newblock  Cambridge University Press (Cambridge Studies in Advanced Mathematics, 2013).	
		
		%\bibitem{M15}
		%P.~Mattila.
		%\newblock {\em Fourier Analysis and Hausdorff Dimension}.
		%\newblock Cambridge Studies in Advanced Mathematics. Cambridge University Press, 2015.

\bibitem{Peral1980}
		J.C.~Peral.
		\newblock $L^p$ estimates for the wave equation.
		\newblock {\em J. of Funct. Anal.}, 36(1):114--145, 1980.
	
		
		\bibitem{SSS91}
		A. Seeger, C.D. Sogge and E.M.~Stein.
		\newblock {\em Regularity Properties of Fourier Integral Operators}.
		\newblock Ann. Math., 134(2): 231--251, 1991.
	

		\bibitem{soggeNLW}
		C.D.~Sogge.
		\newblock {\em Lectures on Non-Linear Wave Equations}.
		\newblock 2nd Edition (International Press of Boston, Inc. 2008).	
		
		\bibitem{S70}
		E.M.~Stein.
		\newblock {\em Singular Integrals and Differentiability Properties of Functions}.
		\newblock Princeton Mathematical Series. Princeton University Press, 1970.
		
		\bibitem{steinha}
		E.M.~Stein and T.S.~Murphy.
		\newblock {\em Harmonic Analysis (PMS-43): Real-Variable Methods, Orthogonality, and Oscillatory Integrals}.
		\newblock Princeton Mathematical Series (Princeton University Press, 1993).
		
		\bibitem{steinshak-complex}
		E.M.~Stein and R.~Shakarchi.
		\newblock {\em Complex Analysis}.
		\newblock Princeton Lectures in Analysis, No. 2. (Princeton University Press, 2003).
		
		\bibitem{steinweiss}
		E.M.~Stein and G.~Weiss.
		\newblock {\em Introduction to Fourier Analysis on Euclidean Spaces}.
		\newblock Princeton Mathematical Series (Princeton University Press, 1971).
		
		
		\bibitem{ST20}
		C.~Sun and N.~Tzvetkov.
		\newblock Concerning the pathological set in the context of probabilistic well-posedness.
		\newblock {\em C. R. Math.}, 358(9-10):989–-999, 2020.	
		
		\bibitem{Tzvetkov2019-notes}
		M.~Gubinelli, P.~E.~Souganidis, and N.~Tzvetkov.
		\newblock In {\em Singular random dynamics ({C}etraro, 2016)}
		
		
		\bibitem{W65}
		S.~Wainger.
		\newblock {\em Special Trigonometric Series in k - dimensions}.
		\newblock Memoirs of the American Mathematical Society, 1965.	
		
		\bibitem{watson}
		G.N.~Watson.
		\newblock {\em A Treatise on the Theory of Bessel Functions}.
		\newblock Cambridge Mathematical Library (Cambridge University Press, 1922).
		
		\bibitem{wolff}
		T.~Wolff.
		\newblock {\em Lectures on Harmonic Analysis}.
		\newblock AMS University Lecture Notes, 2003.
		
	\end{thebibliography}
	
\end{document}
